\documentclass[output=paper,colorlinks,citecolor=brown
% ,hidelinks
% showindex
]{langscibook}
\ChapterDOI{10.5281/zenodo.20285206}

%This is where you put the authors and their affiliations
\author{Doris L. Payne\affiliation{University of Oregon \& SIL Global}}

%Insert your title here
\title{Middle verbs in Maa: Alternating, \textit{media tantum}, and (spurious) look-alikes}
\abstract{Based on semantics, morphophonology, morphosyntax, and some diachronic evidence, this study distinguishes four types of Maa (Maasai) verbs that might initially appear to be middle forms: those which alternate with simple transitive stems, non-alternating middles (sometimes called \textit{media tantum} in the literature), verbs partially assimilated to the \textit{media tantum} category, and spurious ``look-alikes." Compared to alternating middles, a greater proportion of \textit{media tantum} verbs express human related or human experience concepts. The study sets the stage to understand how irregularities in a system may arise. The investigation is based on text and lexicography materials, collected collaboratively with native speakers.}


\IfFileExists{../localcommands.tex}{
   \addbibresource{../localbibliography.bib}
   %%%%%%%%%%%%%%%%%%%%%%%%%%%%%%%%%%%%%%%%%%%%%
%%%%%%%%%% From LSP %%%%%%%%%%%%%%%%%%%%%%%%%
%%%%%%%%%%%%%%%%%%%%%%%%%%%%%%%%%%%%%%%%%%%%%

\usepackage{tabularx,multicol}
\usepackage{url}
\urlstyle{same}

\usepackage{langsci-optional}
\usepackage{langsci-lgr}
\usepackage{langsci-gb4e}

% NOTE THAT THE FOLLOWING PACKAGES ARE BANNED: 
% - pstricks
% - tipa
%%%%%%%%%%%%%%%%%%%%%%%%%%%%%%%%%%%%%%%%%%%%%
%%%%%%%%%% From ACAL LaTeX template %%%%%%%%%
%%%%%%%%%%%%%%%%%%%%%%%%%%%%%%%%%%%%%%%%%%%%%

%For stacking text, used here in autosegmental diagrams
\usepackage{stackengine}

%To combine rows in tables
\usepackage{multirow}

%Gives some extra formatting options, e.g. underlining/strikeout
\usepackage[normalem]{ulem}

%Use package/command below to create a double-spaced document, if you want one. Uncomment BOTH the package and the command (\doublespacing) to create a doublespaced document, or leave them as is to have a single-spaced document.
%\usepackage{setspace}
%\doublespacing 

 

%use for special OT tableaux symbols like bomb and sad face. must be loaded early on because it doesn't play well with some other packages
\usepackage{fourier-orns}

%Basic math symbols 
\usepackage{pifont}
\usepackage{amssymb}

%%%Gives shortcuts for glossing. The use of this package is NOT explained in the Quick Reference Guide, but the documentation is on CTAN for those that are interested. MJKD finds it handy for glossing. (https://ctan.org/pkg/leipzig?lang=en)
\usepackage{leipzig}

%Tables
% \usepackage{caption} %For table captions 

%For OT-style tableaux
\usepackage[notipa]{ot-tableau}

%highlights text with \hl{text}
\usepackage{color, soul} %note that soul sometimes eats Unicode text witouth telling you. 

%Drawing Syntax Trees
\usepackage[linguistics]{forest}

%This specifies some formatting for the forest trees to make them nicer to look at. Unfortunately this was borrowed by Michael Diercks from someone a while ago and he can't remember who. Apologies and if someone lets him know he will update this.
\forestset{
  nice nodes/.style={
    for tree={
      inner sep=0pt,
      fit=band,
    },
  },
  default preamble=nice nodes,
}

%A couple of packages that seemed to prefer being called toward the end of the preamble

%This package provides macros for typesetting SPE-style phonological rules. I've redefined the \oneof command here, so that there the rule includes a right brace. 
\usepackage{phonrule}
\renewcommand{\oneof}[2][c]{\ensuremath{
    ~\left\{
    \begin{tabular}{#1#1}#2\end{tabular}
    \right\}
    }}

%For using Greek letters outside of math mode.
% \usepackage{textgreek}


%Random, lets us use the XeLaTeX logo. Not important to the template at all.
% \usepackage{metalogo}



%%%%%%%%%%%%%%%%%%%%%%%%%%%%%%%%%%%%%%%%%%%%%
%%%%%%%%%% From Smith paper %%%%%%%%%%%%%%%%%
%%%%%%%%%%%%%%%%%%%%%%%%%%%%%%%%%%%%%%%%%%%%%

\usetikzlibrary{positioning}
\usetikzlibrary{trees}
% \usepackage{pstricks} %pstricks is banned. Use tikz instead. 
% \usepackage{pst-node}
%\usepackage{tikz}

%%%%%%%%%%%%%%%%%%%%%%%%%%%%%%%%%%%%%%%%%%%%%
%%%%%%%%%% From Gluckman paper %%%%%%%%%%%%%%
%%%%%%%%%%%%%%%%%%%%%%%%%%%%%%%%%%%%%%%%%%%%%

\usepackage{amsmath}



%%%%%%%%%%%%%%%%%%%%%%%%%%%%%%%%%%%%%%%%%%%%%
%%%%%%%%%% From Kandybowicz paper %%%%%%%%%%%
%%%%%%%%%%%%%%%%%%%%%%%%%%%%%%%%%%%%%%%%%%%%%

%These are in the .tex, but there's nothing in the "Figures" folder so I'm not sure that it's actually doing anything. 
 
% \graphicspath{ {./Figures/} }

%%%%%%%%%%%%%%%%%%%%%%%%%%%%%%%%%%%%%%%%%%%%%
%%%%%%%%%% From Matte paper %%%%%%%%%%%%%%%%%
%%%%%%%%%%%%%%%%%%%%%%%%%%%%%%%%%%%%%%%%%%%%%

%%%%%%%%%%%%%%%%%%%%%%%%%%%%%%%%%%%%%%%%%%%%%
%%%%%%%%%% From Grollemund paper %%%%%%%%%%%%%%
%%%%%%%%%%%%%%%%%%%%%%%%%%%%%%%%%%%%%%%%%%%%%

% \usepackage{lscape}

%%%%%%%%%%%%%%%%%%%%%%%%%%%%%%%%%%%%%%%%%%%%%
%%%%%%%%%% From Burns paper %%%%%%%%%%%%%%
%%%%%%%%%%%%%%%%%%%%%%%%%%%%%%%%%%%%%%%%%%%%%


% \usepackage{url}
% \urlstyle{same}

\usepackage{listings}
\lstset{basicstyle=\ttfamily,tabsize=2,breaklines=true}

%MJKD commented out - these are standard for chapters in edited volumes, I don't think we need them here in the preamble. 

% \usepackage{tabularx,multicol}
% \usepackage{langsci-basic}
% \usepackage{langsci-gb4e}
% \bibliography{localbibliography}
% \newcommand{\orcid}[1]{}
% \pagenumbering{roman}

% \IfFileExists{../localcommands.tex}{%hack to check whether this is being compiled as part of a collection or standalone
%    %%%%%%%%%%%%%%%%%%%%%%%%%%%%%%%%%%%%%%%%%%%%%
%%%%%%%%%% From LSP %%%%%%%%%%%%%%%%%%%%%%%%%
%%%%%%%%%%%%%%%%%%%%%%%%%%%%%%%%%%%%%%%%%%%%%

\usepackage{tabularx,multicol}
\usepackage{url}
\urlstyle{same}

\usepackage{langsci-optional}
\usepackage{langsci-lgr}
\usepackage{langsci-gb4e}

% NOTE THAT THE FOLLOWING PACKAGES ARE BANNED: 
% - pstricks
% - tipa
%%%%%%%%%%%%%%%%%%%%%%%%%%%%%%%%%%%%%%%%%%%%%
%%%%%%%%%% From ACAL LaTeX template %%%%%%%%%
%%%%%%%%%%%%%%%%%%%%%%%%%%%%%%%%%%%%%%%%%%%%%

%For stacking text, used here in autosegmental diagrams
\usepackage{stackengine}

%To combine rows in tables
\usepackage{multirow}

%Gives some extra formatting options, e.g. underlining/strikeout
\usepackage[normalem]{ulem}

%Use package/command below to create a double-spaced document, if you want one. Uncomment BOTH the package and the command (\doublespacing) to create a doublespaced document, or leave them as is to have a single-spaced document.
%\usepackage{setspace}
%\doublespacing 

 

%use for special OT tableaux symbols like bomb and sad face. must be loaded early on because it doesn't play well with some other packages
\usepackage{fourier-orns}

%Basic math symbols 
\usepackage{pifont}
\usepackage{amssymb}

%%%Gives shortcuts for glossing. The use of this package is NOT explained in the Quick Reference Guide, but the documentation is on CTAN for those that are interested. MJKD finds it handy for glossing. (https://ctan.org/pkg/leipzig?lang=en)
\usepackage{leipzig}

%Tables
% \usepackage{caption} %For table captions 

%For OT-style tableaux
\usepackage[notipa]{ot-tableau}

%highlights text with \hl{text}
\usepackage{color, soul} %note that soul sometimes eats Unicode text witouth telling you. 

%Drawing Syntax Trees
\usepackage[linguistics]{forest}

%This specifies some formatting for the forest trees to make them nicer to look at. Unfortunately this was borrowed by Michael Diercks from someone a while ago and he can't remember who. Apologies and if someone lets him know he will update this.
\forestset{
  nice nodes/.style={
    for tree={
      inner sep=0pt,
      fit=band,
    },
  },
  default preamble=nice nodes,
}

%A couple of packages that seemed to prefer being called toward the end of the preamble

%This package provides macros for typesetting SPE-style phonological rules. I've redefined the \oneof command here, so that there the rule includes a right brace. 
\usepackage{phonrule}
\renewcommand{\oneof}[2][c]{\ensuremath{
    ~\left\{
    \begin{tabular}{#1#1}#2\end{tabular}
    \right\}
    }}

%For using Greek letters outside of math mode.
% \usepackage{textgreek}


%Random, lets us use the XeLaTeX logo. Not important to the template at all.
% \usepackage{metalogo}



%%%%%%%%%%%%%%%%%%%%%%%%%%%%%%%%%%%%%%%%%%%%%
%%%%%%%%%% From Smith paper %%%%%%%%%%%%%%%%%
%%%%%%%%%%%%%%%%%%%%%%%%%%%%%%%%%%%%%%%%%%%%%

\usetikzlibrary{positioning}
\usetikzlibrary{trees}
% \usepackage{pstricks} %pstricks is banned. Use tikz instead. 
% \usepackage{pst-node}
%\usepackage{tikz}

%%%%%%%%%%%%%%%%%%%%%%%%%%%%%%%%%%%%%%%%%%%%%
%%%%%%%%%% From Gluckman paper %%%%%%%%%%%%%%
%%%%%%%%%%%%%%%%%%%%%%%%%%%%%%%%%%%%%%%%%%%%%

\usepackage{amsmath}



%%%%%%%%%%%%%%%%%%%%%%%%%%%%%%%%%%%%%%%%%%%%%
%%%%%%%%%% From Kandybowicz paper %%%%%%%%%%%
%%%%%%%%%%%%%%%%%%%%%%%%%%%%%%%%%%%%%%%%%%%%%

%These are in the .tex, but there's nothing in the "Figures" folder so I'm not sure that it's actually doing anything. 
 
% \graphicspath{ {./Figures/} }

%%%%%%%%%%%%%%%%%%%%%%%%%%%%%%%%%%%%%%%%%%%%%
%%%%%%%%%% From Matte paper %%%%%%%%%%%%%%%%%
%%%%%%%%%%%%%%%%%%%%%%%%%%%%%%%%%%%%%%%%%%%%%

%%%%%%%%%%%%%%%%%%%%%%%%%%%%%%%%%%%%%%%%%%%%%
%%%%%%%%%% From Grollemund paper %%%%%%%%%%%%%%
%%%%%%%%%%%%%%%%%%%%%%%%%%%%%%%%%%%%%%%%%%%%%

% \usepackage{lscape}

%%%%%%%%%%%%%%%%%%%%%%%%%%%%%%%%%%%%%%%%%%%%%
%%%%%%%%%% From Burns paper %%%%%%%%%%%%%%
%%%%%%%%%%%%%%%%%%%%%%%%%%%%%%%%%%%%%%%%%%%%%


% \usepackage{url}
% \urlstyle{same}

\usepackage{listings}
\lstset{basicstyle=\ttfamily,tabsize=2,breaklines=true}

%MJKD commented out - these are standard for chapters in edited volumes, I don't think we need them here in the preamble. 

% \usepackage{tabularx,multicol}
% \usepackage{langsci-basic}
% \usepackage{langsci-gb4e}
% \bibliography{localbibliography}
% \newcommand{\orcid}[1]{}
% \pagenumbering{roman}

% \IfFileExists{../localcommands.tex}{%hack to check whether this is being compiled as part of a collection or standalone
%    %%%%%%%%%%%%%%%%%%%%%%%%%%%%%%%%%%%%%%%%%%%%%
%%%%%%%%%% From LSP %%%%%%%%%%%%%%%%%%%%%%%%%
%%%%%%%%%%%%%%%%%%%%%%%%%%%%%%%%%%%%%%%%%%%%%

\usepackage{tabularx,multicol}
\usepackage{url}
\urlstyle{same}

\usepackage{langsci-optional}
\usepackage{langsci-lgr}
\usepackage{langsci-gb4e}

% NOTE THAT THE FOLLOWING PACKAGES ARE BANNED: 
% - pstricks
% - tipa
%%%%%%%%%%%%%%%%%%%%%%%%%%%%%%%%%%%%%%%%%%%%%
%%%%%%%%%% From ACAL LaTeX template %%%%%%%%%
%%%%%%%%%%%%%%%%%%%%%%%%%%%%%%%%%%%%%%%%%%%%%

%For stacking text, used here in autosegmental diagrams
\usepackage{stackengine}

%To combine rows in tables
\usepackage{multirow}

%Gives some extra formatting options, e.g. underlining/strikeout
\usepackage[normalem]{ulem}

%Use package/command below to create a double-spaced document, if you want one. Uncomment BOTH the package and the command (\doublespacing) to create a doublespaced document, or leave them as is to have a single-spaced document.
%\usepackage{setspace}
%\doublespacing 

 

%use for special OT tableaux symbols like bomb and sad face. must be loaded early on because it doesn't play well with some other packages
\usepackage{fourier-orns}

%Basic math symbols 
\usepackage{pifont}
\usepackage{amssymb}

%%%Gives shortcuts for glossing. The use of this package is NOT explained in the Quick Reference Guide, but the documentation is on CTAN for those that are interested. MJKD finds it handy for glossing. (https://ctan.org/pkg/leipzig?lang=en)
\usepackage{leipzig}

%Tables
% \usepackage{caption} %For table captions 

%For OT-style tableaux
\usepackage[notipa]{ot-tableau}

%highlights text with \hl{text}
\usepackage{color, soul} %note that soul sometimes eats Unicode text witouth telling you. 

%Drawing Syntax Trees
\usepackage[linguistics]{forest}

%This specifies some formatting for the forest trees to make them nicer to look at. Unfortunately this was borrowed by Michael Diercks from someone a while ago and he can't remember who. Apologies and if someone lets him know he will update this.
\forestset{
  nice nodes/.style={
    for tree={
      inner sep=0pt,
      fit=band,
    },
  },
  default preamble=nice nodes,
}

%A couple of packages that seemed to prefer being called toward the end of the preamble

%This package provides macros for typesetting SPE-style phonological rules. I've redefined the \oneof command here, so that there the rule includes a right brace. 
\usepackage{phonrule}
\renewcommand{\oneof}[2][c]{\ensuremath{
    ~\left\{
    \begin{tabular}{#1#1}#2\end{tabular}
    \right\}
    }}

%For using Greek letters outside of math mode.
% \usepackage{textgreek}


%Random, lets us use the XeLaTeX logo. Not important to the template at all.
% \usepackage{metalogo}



%%%%%%%%%%%%%%%%%%%%%%%%%%%%%%%%%%%%%%%%%%%%%
%%%%%%%%%% From Smith paper %%%%%%%%%%%%%%%%%
%%%%%%%%%%%%%%%%%%%%%%%%%%%%%%%%%%%%%%%%%%%%%

\usetikzlibrary{positioning}
\usetikzlibrary{trees}
% \usepackage{pstricks} %pstricks is banned. Use tikz instead. 
% \usepackage{pst-node}
%\usepackage{tikz}

%%%%%%%%%%%%%%%%%%%%%%%%%%%%%%%%%%%%%%%%%%%%%
%%%%%%%%%% From Gluckman paper %%%%%%%%%%%%%%
%%%%%%%%%%%%%%%%%%%%%%%%%%%%%%%%%%%%%%%%%%%%%

\usepackage{amsmath}



%%%%%%%%%%%%%%%%%%%%%%%%%%%%%%%%%%%%%%%%%%%%%
%%%%%%%%%% From Kandybowicz paper %%%%%%%%%%%
%%%%%%%%%%%%%%%%%%%%%%%%%%%%%%%%%%%%%%%%%%%%%

%These are in the .tex, but there's nothing in the "Figures" folder so I'm not sure that it's actually doing anything. 
 
% \graphicspath{ {./Figures/} }

%%%%%%%%%%%%%%%%%%%%%%%%%%%%%%%%%%%%%%%%%%%%%
%%%%%%%%%% From Matte paper %%%%%%%%%%%%%%%%%
%%%%%%%%%%%%%%%%%%%%%%%%%%%%%%%%%%%%%%%%%%%%%

%%%%%%%%%%%%%%%%%%%%%%%%%%%%%%%%%%%%%%%%%%%%%
%%%%%%%%%% From Grollemund paper %%%%%%%%%%%%%%
%%%%%%%%%%%%%%%%%%%%%%%%%%%%%%%%%%%%%%%%%%%%%

% \usepackage{lscape}

%%%%%%%%%%%%%%%%%%%%%%%%%%%%%%%%%%%%%%%%%%%%%
%%%%%%%%%% From Burns paper %%%%%%%%%%%%%%
%%%%%%%%%%%%%%%%%%%%%%%%%%%%%%%%%%%%%%%%%%%%%


% \usepackage{url}
% \urlstyle{same}

\usepackage{listings}
\lstset{basicstyle=\ttfamily,tabsize=2,breaklines=true}

%MJKD commented out - these are standard for chapters in edited volumes, I don't think we need them here in the preamble. 

% \usepackage{tabularx,multicol}
% \usepackage{langsci-basic}
% \usepackage{langsci-gb4e}
% \bibliography{localbibliography}
% \newcommand{\orcid}[1]{}
% \pagenumbering{roman}

% \IfFileExists{../localcommands.tex}{%hack to check whether this is being compiled as part of a collection or standalone
%    \input{../localpackages}
%    \input{../localcommands}
%    \input{../localhyphenation}
%     \bibliography{localbibliography}
%     \togglepaper[23]
% }{}

% %custom footer for preprints
% \papernote{\scriptsize\normalfont
%     Change author.
%     Change title.
%     To appear in:
%     Change Volume Editor.  
%     Change volume title.  
%     Berlin: Language Science Press. [preliminary page numbering]
% }



%%%%%%%%%%%%%%%%%%%%%%%%%%%%%%%%%%%%%%%%%%%%%
%%%%%%%%%% From Feibig paper %%%%%%%%%%%%%%
%%%%%%%%%%%%%%%%%%%%%%%%%%%%%%%%%%%%%%%%%%%%%

\usepackage{tikz-qtree}
\usepackage{tikz-qtree-compat}

%%%%%%%%%%%%%%%%%%%%%%%%%%%%%%%%%%%%%%%%%%%%%
%%%%%%%%%% From Olson paper %%%%%%%%%%%%%%
%%%%%%%%%%%%%%%%%%%%%%%%%%%%%%%%%%%%%%%%%%%%%

%MJKD - TIPA is forbidden lol. Commenting out for now but we'll have to address this in the Olson paper

%\usepackage[tone]{tipa}

%%%%%%%%%%%%%%%%%%%%%%%%%%%%%%%%%%%%%%%%%%%%%
%%%%%%%%%% From Caragine paper %%%%%%%%%%%%%%
%%%%%%%%%%%%%%%%%%%%%%%%%%%%%%%%%%%%%%%%%%%%%

\usepackage{placeins}
% \usepackage{graphicx}
% \usepackage{float} %Overleaf suggested this as a solution


%%%%%%%%%%%%%%%%%%%%%%%%%%%%%%%%%%%%%%%%%%%%%
%%%%%%%%%% From Kieffer paper %%%%%%%%%%%%%%
%%%%%%%%%%%%%%%%%%%%%%%%%%%%%%%%%%%%%%%%%%%%%

\usepackage{array}

%%%%%%%%%%%%%%%%%%%%%%%%%%%%%%%%%%%%%%%%%%%%%
%%%%%%%%%% From Payne paper %%%%%%%%%%%%%%
%%%%%%%%%%%%%%%%%%%%%%%%%%%%%%%%%%%%%%%%%%%%%

\usepackage{enumerate}
\usepackage{array,ragged2e}
\usepackage{pgfplots}


%%%%%%%%%%%%%%%%%%%%%%%%%%%%%%%%%%%%%%%%%%%%%
%%%%%%%%%% From Diercks paper %%%%%%%%%%%%%%
%%%%%%%%%%%%%%%%%%%%%%%%%%%%%%%%%%%%%%%%%%%%%

% \usepackage{cgloss}



%%%%%%%%%%%%%%%%%%%%%%%%%%%%%%%%%%%%%%%%%%%%%
%%%%%%%%%% From Li paper %%%%%%%%%%%%%%
%%%%%%%%%%%%%%%%%%%%%%%%%%%%%%%%%%%%%%%%%%%%%



%%%%%%%%%%%%%%%%%%%%%%%%%%%%%%%%%%%%%%%%%%%%%
%%%%%%%%%% From Myers paper %%%%%%%%%%%%%%
%%%%%%%%%%%%%%%%%%%%%%%%%%%%%%%%%%%%%%%%%%%%%



\usepackage{tikz-qtree}
\usepackage{tikz-qtree-compat}


%Package that makes glossing slightly easier.
\RequirePackage{leipzig}
%\RequirePackage{mathtools} % for \mathrlap

% % % Shortcuts (various borrowed from Jason Zentz)
\RequirePackage{xspace}
\xspaceaddexceptions{]\}}
\usepackage{langsci-textipa}
\usepackage{langsci-branding}
\usepackage{tabto}

%    %%%%%%%%%%%%%%%%%%%%%%%%%%%%%%%%%%%%%%%%%%%%%
%%%%%%%%%% From LSP %%%%%%%%%%%%%%%%%%%%%%%%%
%%%%%%%%%%%%%%%%%%%%%%%%%%%%%%%%%%%%%%%%%%%%%


% \newcommand*{\orcid}{}


\makeatletter
\let\thetitle\@title
\let\theauthor\@author
\makeatother

\newcommand{\togglepaper}[1][0]{
   \bibliography{../localbibliography}
   \papernote{\scriptsize\normalfont
     \theauthor.
     \titleTemp.
     To appear in:
     E. Di Tor \& Herr Rausgeberin (ed.).
     Booktitle in localcommands.tex.
     Berlin: Language Science Press. [preliminary page numbering]
   }
   \pagenumbering{roman}
   \setcounter{chapter}{#1}
   \addtocounter{chapter}{-1}
}

\newbool{bookcompile}
\booltrue{bookcompile}
\newcommand{\bookorchapter}[2]{\ifbool{bookcompile}{#1}{#2}}


%%%%%%%%%%%%%%%%%%%%%%%%%%%%%%%%%%%%%%%%%%%%%
%%%%%%%%%% From ACAL LaTeX template %%%%%%%%%
%%%%%%%%%%%%%%%%%%%%%%%%%%%%%%%%%%%%%%%%%%%%%


\usepackage{tikz}

\newcommand{\circled}[1]{\begin{tikzpicture}[baseline=(word.base)]
\node[draw, rounded corners, text height=8pt, text depth=2pt, inner sep=2pt, outer sep=0pt, use as bounding box] (word) {#1};
\end{tikzpicture}
}


%%%%%%%%%%%%%%%%%%%%%%%%%%%%%%%%%%%%%%%%%%%%%%
%%%%%%%%%%%%%%%%%%%%%%%%%%%%%%%%%%%%%%%%%%%%%%

% Useful Ling Shortcuts


%This makes the \emptyset command be a nicer one
\let\oldemptyset\emptyset
\let\emptyset\varnothing
\newcommand{\nothing}{$\emptyset$}

%Not all of these are explained in the Quick Reference Guide, but they are here bc they are relevant to some of our students. 
\newcommand{\xbar}{\ensuremath{'}\xspace}
\newcommand{\xhead}{\ensuremath{^\circ}\xspace}
\newcommand{\Lb}[1]{$\text{[}_{\text{#1}}$ } %A more convenient left bracket
\newcommand{\Rb}[1]{$\text{]}_{\text{#1}}$ } %A more convenient left bracket
\newcommand{\gap}{\underline{\hspace{1.2em}}}
\newcommand{\vP}{\emph{v}P}
\newcommand{\lilv}{\emph{v}}
\newcommand{\Abar}{A$'$-} %A more convenient A-bar notation
\newcommand{\ph}{$\varphi$\xspace} %A more convenient phi
\newcommand{\pro}{\emph{pro}\xspace}
\newcommand{\subs}[1]{\textsubscript{#1}} %A more convenient subscript
\newcommand{\spells}{$\Longleftrightarrow$} %spellout arrow for morph spellout rules
\newcommand{\tr}[1]{\textit{t}\textsubscript{\textit{#1}}} %easy traces with subscript
\newcommand{\supers}[1]{\textsuperscript{#1}}

%These are borrowed/modified from Andrew Mackenzie's ling-macros package. We have copied them in here (instead of just using the package itself) to avoid a minor clash that was generating an error.


% Abbreviations for glossing, based on Leipzig
\newleipzig{hab}{hab}{habitual}
\newleipzig{rem}{rem}{remote}
\newleipzig{sm}{sm}{subject marker}
\newleipzig{t}{t}{tense}
\newleipzig{aa}{aa}{anti-agreement}
\newleipzig{pron}{pron}{pronoun}
\newleipzig{rec}{rec}{recent}
\newleipzig{om}{om}{object marker}
%\newleipzig{ipfv}{ipfv}{imperfective}
\newleipzig{asp}{asp}{aspect}
\newleipzig{lk}{lk}{linker}
\newleipzig{pcl}{pcl}{particle}
\newleipzig{stat}{stat}{stative}
\newleipzig{ints}{ints}{intensive}
\newleipzig{ascl}{ascl}{assertive subject clitic}
\newleipzig{nascl}{nascl}{non-assertive subject clitic}
\newleipzig{ta}{ta}{tense and/or aspect}
\newleipzig{assoc}{assoc}{associative marker}
\newleipzig{hon}{hon}{honorific}
%\newleipzig{whprt}{wh}{\wh particle}
\newleipzig{sa}{sa}{subject agreement}
\newleipzig{conj}{conj}{conjunction}
%\newleipzig{loc}{loc}{locative}
\newleipzig{expl}{expl}{expletive}
\newleipzig{rcm}{rcm}{reciprocal marker}
\newleipzig{pers}{pers}{persistive}
\newleipzig{fv}{fv}{final vowel}
%\newleipzig{}{}{} %this is just to copy for when I want to add more


%%%%%%%%%%%%%%%%%%%%%%%%%%%%%%%%%%%%%%%%%%%%%%%%%%%%%
%%%%%%%%%%%%%%%%%%%%%%%%%%%%%%%%%%%%%%%%%%%%%%%%%%%%%


%%%%%%%%%%%%%%%%%%%%%%%%%%%%%%%%%%%%%%%%%%%%%
%%%%%%%%%% From Gluckman paper %%%%%%%%%%%%%%
%%%%%%%%%%%%%%%%%%%%%%%%%%%%%%%%%%%%%%%%%%%%%

\newcommand{\pcite}[1]{\citeauthor{#1}'s (\citeyear{#1})}


%%%%%%%%%%%%%%%%%%%%%%%%%%%%%%%%%%%%%%%%%%%%%
%%%%%%%%%% From Myers paper %%%%%%%%%%%%%%
%%%%%%%%%%%%%%%%%%%%%%%%%%%%%%%%%%%%%%%%%%%%%

%% hspace over width of something without showing it
\newcommand*{\hspaceThis}[1]{\settowidth{\LSPTmp}{#1}\hspace*{\LSPTmp}}
% use \hphantom{} for this


%%%%%%%%%%%%%%%%%%%%%%%%%%%%%%%%%%%%%%%%%%%%%
%%%%%%%%%% From Matte paper %%%%%%%%%%%%%%%%%
%%%%%%%%%%%%%%%%%%%%%%%%%%%%%%%%%%%%%%%%%%%%%

%mjkd commented this out in case it's breaking something right now. 
% \newcounter{xlistex}
% \newcommand{\footnoteex}{\stepcounter{xlistex}(\roman{xlistex})}

\newleipzig{sp}{sp}{speaker} %a new leipzig glossing command


%%%%%%%%%%%%%%%%%%%%%%%%%%%%%%%%%%%%%%%%%%%%%
%%%%%%%%%% From GamFranich paper %%%%%%%%%%%%%%
%%%%%%%%%%%%%%%%%%%%%%%%%%%%%%%%%%%%%%%%%%%%%

%MJKD commented these out - both will be provided by the overall LSP book template 

% \bibliography{localbibliography}
% \newcommand{\orcid}[1]{}

%%%%%%%%%%%%%%%%%%%%%%%%%%%%%%%%%%%%%%%%%%%%%
%%%%%%%%%% From Diercks paper %%%%%%%%%%%%%%
%%%%%%%%%%%%%%%%%%%%%%%%%%%%%%%%%%%%%%%%%%%%%

\newcommand{\abar}{$\bar{\text{A}}$\xspace} % this one is different to the \Abar command in Mike's shortcuts doc
\newcommand{\topFeat}{[\textsc{top}]\xspace}


%%%%%%%%%%%%%%%%%%%%%%%%%%%%%%%%%%%%%%%%%%%%%
%%%%%%%%%% From Kinchsular paper %%%%%%%%%%%%%%
%%%%%%%%%%%%%%%%%%%%%%%%%%%%%%%%%%%%%%%%%%%%%

%%%%%%%%%%%%%%%%%%%%%%%%%%%%%%%%%%%%%%%%%%%%%
%%%%%%%%%% From Chen paper %%%%%%%%%%%%%%
%%%%%%%%%%%%%%%%%%%%%%%%%%%%%%%%%%%%%%%%%%%%%

\newcounter{savedex}
\makeatletter
\newcommand*{\saveex}{\setcounter{savedex}{\value{\@xnumctr}}}
\newcommand*{\resumeex}{\setcounter{\@xnumctr}{\value{savedex}}}
\makeatother


%    \input{../localhyphenation}
%     \bibliography{localbibliography}
%     \togglepaper[23]
% }{}

% %custom footer for preprints
% \papernote{\scriptsize\normalfont
%     Change author.
%     Change title.
%     To appear in:
%     Change Volume Editor.  
%     Change volume title.  
%     Berlin: Language Science Press. [preliminary page numbering]
% }



%%%%%%%%%%%%%%%%%%%%%%%%%%%%%%%%%%%%%%%%%%%%%
%%%%%%%%%% From Feibig paper %%%%%%%%%%%%%%
%%%%%%%%%%%%%%%%%%%%%%%%%%%%%%%%%%%%%%%%%%%%%

\usepackage{tikz-qtree}
\usepackage{tikz-qtree-compat}

%%%%%%%%%%%%%%%%%%%%%%%%%%%%%%%%%%%%%%%%%%%%%
%%%%%%%%%% From Olson paper %%%%%%%%%%%%%%
%%%%%%%%%%%%%%%%%%%%%%%%%%%%%%%%%%%%%%%%%%%%%

%MJKD - TIPA is forbidden lol. Commenting out for now but we'll have to address this in the Olson paper

%\usepackage[tone]{tipa}

%%%%%%%%%%%%%%%%%%%%%%%%%%%%%%%%%%%%%%%%%%%%%
%%%%%%%%%% From Caragine paper %%%%%%%%%%%%%%
%%%%%%%%%%%%%%%%%%%%%%%%%%%%%%%%%%%%%%%%%%%%%

\usepackage{placeins}
% \usepackage{graphicx}
% \usepackage{float} %Overleaf suggested this as a solution


%%%%%%%%%%%%%%%%%%%%%%%%%%%%%%%%%%%%%%%%%%%%%
%%%%%%%%%% From Kieffer paper %%%%%%%%%%%%%%
%%%%%%%%%%%%%%%%%%%%%%%%%%%%%%%%%%%%%%%%%%%%%

\usepackage{array}

%%%%%%%%%%%%%%%%%%%%%%%%%%%%%%%%%%%%%%%%%%%%%
%%%%%%%%%% From Payne paper %%%%%%%%%%%%%%
%%%%%%%%%%%%%%%%%%%%%%%%%%%%%%%%%%%%%%%%%%%%%

\usepackage{enumerate}
\usepackage{array,ragged2e}
\usepackage{pgfplots}


%%%%%%%%%%%%%%%%%%%%%%%%%%%%%%%%%%%%%%%%%%%%%
%%%%%%%%%% From Diercks paper %%%%%%%%%%%%%%
%%%%%%%%%%%%%%%%%%%%%%%%%%%%%%%%%%%%%%%%%%%%%

% \usepackage{cgloss}



%%%%%%%%%%%%%%%%%%%%%%%%%%%%%%%%%%%%%%%%%%%%%
%%%%%%%%%% From Li paper %%%%%%%%%%%%%%
%%%%%%%%%%%%%%%%%%%%%%%%%%%%%%%%%%%%%%%%%%%%%



%%%%%%%%%%%%%%%%%%%%%%%%%%%%%%%%%%%%%%%%%%%%%
%%%%%%%%%% From Myers paper %%%%%%%%%%%%%%
%%%%%%%%%%%%%%%%%%%%%%%%%%%%%%%%%%%%%%%%%%%%%



\usepackage{tikz-qtree}
\usepackage{tikz-qtree-compat}


%Package that makes glossing slightly easier.
\RequirePackage{leipzig}
%\RequirePackage{mathtools} % for \mathrlap

% % % Shortcuts (various borrowed from Jason Zentz)
\RequirePackage{xspace}
\xspaceaddexceptions{]\}}
\usepackage{langsci-textipa}
\usepackage{langsci-branding}
\usepackage{tabto}

%    %%%%%%%%%%%%%%%%%%%%%%%%%%%%%%%%%%%%%%%%%%%%%
%%%%%%%%%% From LSP %%%%%%%%%%%%%%%%%%%%%%%%%
%%%%%%%%%%%%%%%%%%%%%%%%%%%%%%%%%%%%%%%%%%%%%


% \newcommand*{\orcid}{}


\makeatletter
\let\thetitle\@title
\let\theauthor\@author
\makeatother

\newcommand{\togglepaper}[1][0]{
   \bibliography{../localbibliography}
   \papernote{\scriptsize\normalfont
     \theauthor.
     \titleTemp.
     To appear in:
     E. Di Tor \& Herr Rausgeberin (ed.).
     Booktitle in localcommands.tex.
     Berlin: Language Science Press. [preliminary page numbering]
   }
   \pagenumbering{roman}
   \setcounter{chapter}{#1}
   \addtocounter{chapter}{-1}
}

\newbool{bookcompile}
\booltrue{bookcompile}
\newcommand{\bookorchapter}[2]{\ifbool{bookcompile}{#1}{#2}}


%%%%%%%%%%%%%%%%%%%%%%%%%%%%%%%%%%%%%%%%%%%%%
%%%%%%%%%% From ACAL LaTeX template %%%%%%%%%
%%%%%%%%%%%%%%%%%%%%%%%%%%%%%%%%%%%%%%%%%%%%%


\usepackage{tikz}

\newcommand{\circled}[1]{\begin{tikzpicture}[baseline=(word.base)]
\node[draw, rounded corners, text height=8pt, text depth=2pt, inner sep=2pt, outer sep=0pt, use as bounding box] (word) {#1};
\end{tikzpicture}
}


%%%%%%%%%%%%%%%%%%%%%%%%%%%%%%%%%%%%%%%%%%%%%%
%%%%%%%%%%%%%%%%%%%%%%%%%%%%%%%%%%%%%%%%%%%%%%

% Useful Ling Shortcuts


%This makes the \emptyset command be a nicer one
\let\oldemptyset\emptyset
\let\emptyset\varnothing
\newcommand{\nothing}{$\emptyset$}

%Not all of these are explained in the Quick Reference Guide, but they are here bc they are relevant to some of our students. 
\newcommand{\xbar}{\ensuremath{'}\xspace}
\newcommand{\xhead}{\ensuremath{^\circ}\xspace}
\newcommand{\Lb}[1]{$\text{[}_{\text{#1}}$ } %A more convenient left bracket
\newcommand{\Rb}[1]{$\text{]}_{\text{#1}}$ } %A more convenient left bracket
\newcommand{\gap}{\underline{\hspace{1.2em}}}
\newcommand{\vP}{\emph{v}P}
\newcommand{\lilv}{\emph{v}}
\newcommand{\Abar}{A$'$-} %A more convenient A-bar notation
\newcommand{\ph}{$\varphi$\xspace} %A more convenient phi
\newcommand{\pro}{\emph{pro}\xspace}
\newcommand{\subs}[1]{\textsubscript{#1}} %A more convenient subscript
\newcommand{\spells}{$\Longleftrightarrow$} %spellout arrow for morph spellout rules
\newcommand{\tr}[1]{\textit{t}\textsubscript{\textit{#1}}} %easy traces with subscript
\newcommand{\supers}[1]{\textsuperscript{#1}}

%These are borrowed/modified from Andrew Mackenzie's ling-macros package. We have copied them in here (instead of just using the package itself) to avoid a minor clash that was generating an error.


% Abbreviations for glossing, based on Leipzig
\newleipzig{hab}{hab}{habitual}
\newleipzig{rem}{rem}{remote}
\newleipzig{sm}{sm}{subject marker}
\newleipzig{t}{t}{tense}
\newleipzig{aa}{aa}{anti-agreement}
\newleipzig{pron}{pron}{pronoun}
\newleipzig{rec}{rec}{recent}
\newleipzig{om}{om}{object marker}
%\newleipzig{ipfv}{ipfv}{imperfective}
\newleipzig{asp}{asp}{aspect}
\newleipzig{lk}{lk}{linker}
\newleipzig{pcl}{pcl}{particle}
\newleipzig{stat}{stat}{stative}
\newleipzig{ints}{ints}{intensive}
\newleipzig{ascl}{ascl}{assertive subject clitic}
\newleipzig{nascl}{nascl}{non-assertive subject clitic}
\newleipzig{ta}{ta}{tense and/or aspect}
\newleipzig{assoc}{assoc}{associative marker}
\newleipzig{hon}{hon}{honorific}
%\newleipzig{whprt}{wh}{\wh particle}
\newleipzig{sa}{sa}{subject agreement}
\newleipzig{conj}{conj}{conjunction}
%\newleipzig{loc}{loc}{locative}
\newleipzig{expl}{expl}{expletive}
\newleipzig{rcm}{rcm}{reciprocal marker}
\newleipzig{pers}{pers}{persistive}
\newleipzig{fv}{fv}{final vowel}
%\newleipzig{}{}{} %this is just to copy for when I want to add more


%%%%%%%%%%%%%%%%%%%%%%%%%%%%%%%%%%%%%%%%%%%%%%%%%%%%%
%%%%%%%%%%%%%%%%%%%%%%%%%%%%%%%%%%%%%%%%%%%%%%%%%%%%%


%%%%%%%%%%%%%%%%%%%%%%%%%%%%%%%%%%%%%%%%%%%%%
%%%%%%%%%% From Gluckman paper %%%%%%%%%%%%%%
%%%%%%%%%%%%%%%%%%%%%%%%%%%%%%%%%%%%%%%%%%%%%

\newcommand{\pcite}[1]{\citeauthor{#1}'s (\citeyear{#1})}


%%%%%%%%%%%%%%%%%%%%%%%%%%%%%%%%%%%%%%%%%%%%%
%%%%%%%%%% From Myers paper %%%%%%%%%%%%%%
%%%%%%%%%%%%%%%%%%%%%%%%%%%%%%%%%%%%%%%%%%%%%

%% hspace over width of something without showing it
\newcommand*{\hspaceThis}[1]{\settowidth{\LSPTmp}{#1}\hspace*{\LSPTmp}}
% use \hphantom{} for this


%%%%%%%%%%%%%%%%%%%%%%%%%%%%%%%%%%%%%%%%%%%%%
%%%%%%%%%% From Matte paper %%%%%%%%%%%%%%%%%
%%%%%%%%%%%%%%%%%%%%%%%%%%%%%%%%%%%%%%%%%%%%%

%mjkd commented this out in case it's breaking something right now. 
% \newcounter{xlistex}
% \newcommand{\footnoteex}{\stepcounter{xlistex}(\roman{xlistex})}

\newleipzig{sp}{sp}{speaker} %a new leipzig glossing command


%%%%%%%%%%%%%%%%%%%%%%%%%%%%%%%%%%%%%%%%%%%%%
%%%%%%%%%% From GamFranich paper %%%%%%%%%%%%%%
%%%%%%%%%%%%%%%%%%%%%%%%%%%%%%%%%%%%%%%%%%%%%

%MJKD commented these out - both will be provided by the overall LSP book template 

% \bibliography{localbibliography}
% \newcommand{\orcid}[1]{}

%%%%%%%%%%%%%%%%%%%%%%%%%%%%%%%%%%%%%%%%%%%%%
%%%%%%%%%% From Diercks paper %%%%%%%%%%%%%%
%%%%%%%%%%%%%%%%%%%%%%%%%%%%%%%%%%%%%%%%%%%%%

\newcommand{\abar}{$\bar{\text{A}}$\xspace} % this one is different to the \Abar command in Mike's shortcuts doc
\newcommand{\topFeat}{[\textsc{top}]\xspace}


%%%%%%%%%%%%%%%%%%%%%%%%%%%%%%%%%%%%%%%%%%%%%
%%%%%%%%%% From Kinchsular paper %%%%%%%%%%%%%%
%%%%%%%%%%%%%%%%%%%%%%%%%%%%%%%%%%%%%%%%%%%%%

%%%%%%%%%%%%%%%%%%%%%%%%%%%%%%%%%%%%%%%%%%%%%
%%%%%%%%%% From Chen paper %%%%%%%%%%%%%%
%%%%%%%%%%%%%%%%%%%%%%%%%%%%%%%%%%%%%%%%%%%%%

\newcounter{savedex}
\makeatletter
\newcommand*{\saveex}{\setcounter{savedex}{\value{\@xnumctr}}}
\newcommand*{\resumeex}{\setcounter{\@xnumctr}{\value{savedex}}}
\makeatother


%    \input{../localhyphenation}
%     \bibliography{localbibliography}
%     \togglepaper[23]
% }{}

% %custom footer for preprints
% \papernote{\scriptsize\normalfont
%     Change author.
%     Change title.
%     To appear in:
%     Change Volume Editor.  
%     Change volume title.  
%     Berlin: Language Science Press. [preliminary page numbering]
% }



%%%%%%%%%%%%%%%%%%%%%%%%%%%%%%%%%%%%%%%%%%%%%
%%%%%%%%%% From Feibig paper %%%%%%%%%%%%%%
%%%%%%%%%%%%%%%%%%%%%%%%%%%%%%%%%%%%%%%%%%%%%

\usepackage{tikz-qtree}
\usepackage{tikz-qtree-compat}

%%%%%%%%%%%%%%%%%%%%%%%%%%%%%%%%%%%%%%%%%%%%%
%%%%%%%%%% From Olson paper %%%%%%%%%%%%%%
%%%%%%%%%%%%%%%%%%%%%%%%%%%%%%%%%%%%%%%%%%%%%

%MJKD - TIPA is forbidden lol. Commenting out for now but we'll have to address this in the Olson paper

%\usepackage[tone]{tipa}

%%%%%%%%%%%%%%%%%%%%%%%%%%%%%%%%%%%%%%%%%%%%%
%%%%%%%%%% From Caragine paper %%%%%%%%%%%%%%
%%%%%%%%%%%%%%%%%%%%%%%%%%%%%%%%%%%%%%%%%%%%%

\usepackage{placeins}
% \usepackage{graphicx}
% \usepackage{float} %Overleaf suggested this as a solution


%%%%%%%%%%%%%%%%%%%%%%%%%%%%%%%%%%%%%%%%%%%%%
%%%%%%%%%% From Kieffer paper %%%%%%%%%%%%%%
%%%%%%%%%%%%%%%%%%%%%%%%%%%%%%%%%%%%%%%%%%%%%

\usepackage{array}

%%%%%%%%%%%%%%%%%%%%%%%%%%%%%%%%%%%%%%%%%%%%%
%%%%%%%%%% From Payne paper %%%%%%%%%%%%%%
%%%%%%%%%%%%%%%%%%%%%%%%%%%%%%%%%%%%%%%%%%%%%

\usepackage{enumerate}
\usepackage{array,ragged2e}
\usepackage{pgfplots}


%%%%%%%%%%%%%%%%%%%%%%%%%%%%%%%%%%%%%%%%%%%%%
%%%%%%%%%% From Diercks paper %%%%%%%%%%%%%%
%%%%%%%%%%%%%%%%%%%%%%%%%%%%%%%%%%%%%%%%%%%%%

% \usepackage{cgloss}



%%%%%%%%%%%%%%%%%%%%%%%%%%%%%%%%%%%%%%%%%%%%%
%%%%%%%%%% From Li paper %%%%%%%%%%%%%%
%%%%%%%%%%%%%%%%%%%%%%%%%%%%%%%%%%%%%%%%%%%%%



%%%%%%%%%%%%%%%%%%%%%%%%%%%%%%%%%%%%%%%%%%%%%
%%%%%%%%%% From Myers paper %%%%%%%%%%%%%%
%%%%%%%%%%%%%%%%%%%%%%%%%%%%%%%%%%%%%%%%%%%%%



\usepackage{tikz-qtree}
\usepackage{tikz-qtree-compat}


%Package that makes glossing slightly easier.
\RequirePackage{leipzig}
%\RequirePackage{mathtools} % for \mathrlap

% % % Shortcuts (various borrowed from Jason Zentz)
\RequirePackage{xspace}
\xspaceaddexceptions{]\}}
\usepackage{langsci-textipa}
\usepackage{langsci-branding}
\usepackage{tabto}

   %%%%%%%%%%%%%%%%%%%%%%%%%%%%%%%%%%%%%%%%%%%%%
%%%%%%%%%% From LSP %%%%%%%%%%%%%%%%%%%%%%%%%
%%%%%%%%%%%%%%%%%%%%%%%%%%%%%%%%%%%%%%%%%%%%%


% \newcommand*{\orcid}{}


\makeatletter
\let\thetitle\@title
\let\theauthor\@author
\makeatother

\newcommand{\togglepaper}[1][0]{
   \bibliography{../localbibliography}
   \papernote{\scriptsize\normalfont
     \theauthor.
     \titleTemp.
     To appear in:
     E. Di Tor \& Herr Rausgeberin (ed.).
     Booktitle in localcommands.tex.
     Berlin: Language Science Press. [preliminary page numbering]
   }
   \pagenumbering{roman}
   \setcounter{chapter}{#1}
   \addtocounter{chapter}{-1}
}

\newbool{bookcompile}
\booltrue{bookcompile}
\newcommand{\bookorchapter}[2]{\ifbool{bookcompile}{#1}{#2}}


%%%%%%%%%%%%%%%%%%%%%%%%%%%%%%%%%%%%%%%%%%%%%
%%%%%%%%%% From ACAL LaTeX template %%%%%%%%%
%%%%%%%%%%%%%%%%%%%%%%%%%%%%%%%%%%%%%%%%%%%%%


\usepackage{tikz}

\newcommand{\circled}[1]{\begin{tikzpicture}[baseline=(word.base)]
\node[draw, rounded corners, text height=8pt, text depth=2pt, inner sep=2pt, outer sep=0pt, use as bounding box] (word) {#1};
\end{tikzpicture}
}


%%%%%%%%%%%%%%%%%%%%%%%%%%%%%%%%%%%%%%%%%%%%%%
%%%%%%%%%%%%%%%%%%%%%%%%%%%%%%%%%%%%%%%%%%%%%%

% Useful Ling Shortcuts


%This makes the \emptyset command be a nicer one
\let\oldemptyset\emptyset
\let\emptyset\varnothing
\newcommand{\nothing}{$\emptyset$}

%Not all of these are explained in the Quick Reference Guide, but they are here bc they are relevant to some of our students. 
\newcommand{\xbar}{\ensuremath{'}\xspace}
\newcommand{\xhead}{\ensuremath{^\circ}\xspace}
\newcommand{\Lb}[1]{$\text{[}_{\text{#1}}$ } %A more convenient left bracket
\newcommand{\Rb}[1]{$\text{]}_{\text{#1}}$ } %A more convenient left bracket
\newcommand{\gap}{\underline{\hspace{1.2em}}}
\newcommand{\vP}{\emph{v}P}
\newcommand{\lilv}{\emph{v}}
\newcommand{\Abar}{A$'$-} %A more convenient A-bar notation
\newcommand{\ph}{$\varphi$\xspace} %A more convenient phi
\newcommand{\pro}{\emph{pro}\xspace}
\newcommand{\subs}[1]{\textsubscript{#1}} %A more convenient subscript
\newcommand{\spells}{$\Longleftrightarrow$} %spellout arrow for morph spellout rules
\newcommand{\tr}[1]{\textit{t}\textsubscript{\textit{#1}}} %easy traces with subscript
\newcommand{\supers}[1]{\textsuperscript{#1}}

%These are borrowed/modified from Andrew Mackenzie's ling-macros package. We have copied them in here (instead of just using the package itself) to avoid a minor clash that was generating an error.


% Abbreviations for glossing, based on Leipzig
\newleipzig{hab}{hab}{habitual}
\newleipzig{rem}{rem}{remote}
\newleipzig{sm}{sm}{subject marker}
\newleipzig{t}{t}{tense}
\newleipzig{aa}{aa}{anti-agreement}
\newleipzig{pron}{pron}{pronoun}
\newleipzig{rec}{rec}{recent}
\newleipzig{om}{om}{object marker}
%\newleipzig{ipfv}{ipfv}{imperfective}
\newleipzig{asp}{asp}{aspect}
\newleipzig{lk}{lk}{linker}
\newleipzig{pcl}{pcl}{particle}
\newleipzig{stat}{stat}{stative}
\newleipzig{ints}{ints}{intensive}
\newleipzig{ascl}{ascl}{assertive subject clitic}
\newleipzig{nascl}{nascl}{non-assertive subject clitic}
\newleipzig{ta}{ta}{tense and/or aspect}
\newleipzig{assoc}{assoc}{associative marker}
\newleipzig{hon}{hon}{honorific}
%\newleipzig{whprt}{wh}{\wh particle}
\newleipzig{sa}{sa}{subject agreement}
\newleipzig{conj}{conj}{conjunction}
%\newleipzig{loc}{loc}{locative}
\newleipzig{expl}{expl}{expletive}
\newleipzig{rcm}{rcm}{reciprocal marker}
\newleipzig{pers}{pers}{persistive}
\newleipzig{fv}{fv}{final vowel}
%\newleipzig{}{}{} %this is just to copy for when I want to add more


%%%%%%%%%%%%%%%%%%%%%%%%%%%%%%%%%%%%%%%%%%%%%%%%%%%%%
%%%%%%%%%%%%%%%%%%%%%%%%%%%%%%%%%%%%%%%%%%%%%%%%%%%%%


%%%%%%%%%%%%%%%%%%%%%%%%%%%%%%%%%%%%%%%%%%%%%
%%%%%%%%%% From Gluckman paper %%%%%%%%%%%%%%
%%%%%%%%%%%%%%%%%%%%%%%%%%%%%%%%%%%%%%%%%%%%%

\newcommand{\pcite}[1]{\citeauthor{#1}'s (\citeyear{#1})}


%%%%%%%%%%%%%%%%%%%%%%%%%%%%%%%%%%%%%%%%%%%%%
%%%%%%%%%% From Myers paper %%%%%%%%%%%%%%
%%%%%%%%%%%%%%%%%%%%%%%%%%%%%%%%%%%%%%%%%%%%%

%% hspace over width of something without showing it
\newcommand*{\hspaceThis}[1]{\settowidth{\LSPTmp}{#1}\hspace*{\LSPTmp}}
% use \hphantom{} for this


%%%%%%%%%%%%%%%%%%%%%%%%%%%%%%%%%%%%%%%%%%%%%
%%%%%%%%%% From Matte paper %%%%%%%%%%%%%%%%%
%%%%%%%%%%%%%%%%%%%%%%%%%%%%%%%%%%%%%%%%%%%%%

%mjkd commented this out in case it's breaking something right now. 
% \newcounter{xlistex}
% \newcommand{\footnoteex}{\stepcounter{xlistex}(\roman{xlistex})}

\newleipzig{sp}{sp}{speaker} %a new leipzig glossing command


%%%%%%%%%%%%%%%%%%%%%%%%%%%%%%%%%%%%%%%%%%%%%
%%%%%%%%%% From GamFranich paper %%%%%%%%%%%%%%
%%%%%%%%%%%%%%%%%%%%%%%%%%%%%%%%%%%%%%%%%%%%%

%MJKD commented these out - both will be provided by the overall LSP book template 

% \bibliography{localbibliography}
% \newcommand{\orcid}[1]{}

%%%%%%%%%%%%%%%%%%%%%%%%%%%%%%%%%%%%%%%%%%%%%
%%%%%%%%%% From Diercks paper %%%%%%%%%%%%%%
%%%%%%%%%%%%%%%%%%%%%%%%%%%%%%%%%%%%%%%%%%%%%

\newcommand{\abar}{$\bar{\text{A}}$\xspace} % this one is different to the \Abar command in Mike's shortcuts doc
\newcommand{\topFeat}{[\textsc{top}]\xspace}


%%%%%%%%%%%%%%%%%%%%%%%%%%%%%%%%%%%%%%%%%%%%%
%%%%%%%%%% From Kinchsular paper %%%%%%%%%%%%%%
%%%%%%%%%%%%%%%%%%%%%%%%%%%%%%%%%%%%%%%%%%%%%

%%%%%%%%%%%%%%%%%%%%%%%%%%%%%%%%%%%%%%%%%%%%%
%%%%%%%%%% From Chen paper %%%%%%%%%%%%%%
%%%%%%%%%%%%%%%%%%%%%%%%%%%%%%%%%%%%%%%%%%%%%

\newcounter{savedex}
\makeatletter
\newcommand*{\saveex}{\setcounter{savedex}{\value{\@xnumctr}}}
\newcommand*{\resumeex}{\setcounter{\@xnumctr}{\value{savedex}}}
\makeatother


   \input{../localhyphenation}
   \boolfalse{bookcompile}
   \togglepaper[21]%%chapternumber
}{}

\begin{document}
\maketitle

%This just introduces the table of contents for this quick reference guide, and can be deleted if the table of contents is not desired
%\begingroup
%\let\cleardoublepage\clearpage
%\tableofcontents
%\endgroup
\section{The problem: Maa verbs ending in \textit{a} and \textit{o}}
Maa (Eastern Nilotic of Kenya and Tanzania; also known in the literature as \textit{Masai} and \textit{Maasai}; ISO mas) has a set of verbs or verb stems which \citet[134-139]{TuckerMpaayei1955} refer to as ``neuter" forms. These verbs can or sometimes must end in \textit{a} or \textit{o}, and have reflexive/reciprocal, non-active, or in some cases intransitive translational movement meaning. \citet[134]{TuckerMpaayei1955} analyzed \textit{-a/-o} as a ``present neuter" suffix, stating that it alternated with \textit{-ayu/-oyu} in the ``future", and \textit{-ɛ/-e} in the ``past." From a typological perspective, this paper argues that verbs which, at least at first glance, fit the relevant meaning types and end with \textit{a}/\textit{o} do not form a homogeneous group. Most do fit the profile of what are called ``middle" verbs in other linguistic literature. However, this study shows that there are four subgroups of apparent middle verbs, based on semantics, morphophonology, morphosyntax, and some diachronic evidence:

\begin{itemize}
\item Alternating or ``oppositional" \citep{Grestenberger2018, Inglese2022b, Inglese2022a} middles: These productively alternate with simple transitive stems.  
\item Maa-style \textit{media tantum}: These are non-alternating or ``non-oppositional" stems, which in citation form can only occur as a middle.
\item Apparent middle stems which may be partially assimilated to the \textit{media tantum} category.
\item Spurious, though phonological and semantic look-alikes to \textit{media tantum} verbs.
\end{itemize}

The study thus elaborates on Tucker \& Mpaayei's description, and sets the stage to help us better understand the historical development of middle systems and how irregularities may arise. The investigation is, however, not primarily an historical one. 

Data were culled from a list of about 540 verbs which initially appeared to be middles, drawn from lexicography and text material collected collaboratively with native speakers. The paper is organized as follows: \sectref{sec:Payne:Definition} gives a brief treatment of typological approaches to defining ``middle." \sectref{sec:Payne:MaaBackground} introduces basic information on Maa and relevant morphology. \sectref{sec:Payne:Alternating}--\sectref{sec:Payne:Spurious} describe the four subgroups of verbs listed above. \sectref{sec:Payne:Motivation} undertakes a descriptive statistical study of the semantic profiles of alternating versus \textit{media tantum} verbs, in an effort to understand motivations for why \textit{media tantum} verb stems might arise. \sectref{sec:Payne:Conclusions} draws brief conclusions about why categories like ``middle verbs," or others, might ``leak."

\section{Definition of ``middle"}\label{sec:Payne:Definition}

Three characterizations of ``middle" will figure in our discussion of Maa.\footnote{To maintain a coherent focus, I do not address formalist approaches to characterizing middle verbs; see \citet{Embick2004} and \citet{Grestenberger2018} for some introduction to that literature.} \citet[373]{Lyons1968} defines the middle largely semantically as a construction where ``the `action’ or `state’ affects the subject of the verb or his interests.” In contrast, in her typological study of about 32 languages, \citet[3]{Kemmer1993} states that ``At present there is no generally accepted definition or characterization of the middle voice." That work goes on to argue for a general definition of ``middle voice" as involving situation types with a relatively low degree of event elaboration; that is, there is low conceptual difference between the Initiator and the Endpoint of the situation. Kemmer's discussion can be read as positing the reflexive situation as the prototype of the middle category. 

In more recent investigation of the nature and polyfunctionalality of middle voice systems, \citet[490]{Inglese2022b} writes, “The greatest obstacle in the typological study of the middle is that there is no agreement on what cross-linguistically counts as a middle marker.” Inglese finds Kemmer's and others' definitions still lacking for typological work and thus he proposes a multi-part definition of a ``middle system" as involving a construction which:\footnote{The parenthetical comments are my attempt to understand why \citet{Inglese2022a} includes the criterion.} 
\begin{enumerate}[(i)]
\item Occurs with multivalent verbs to encode one or more of the following valency reducing operations: \textsc{passive, reflexive, reciprocal, antipassive, anticausative}. (This excludes patient-oriented argument marking in split-intransitive systems from being considered middle morphology.) 
\item Is obligatory with some monovalent verbs; that is, they cannot occur without a ``middle marker.” (This ensures that some stems are lexicalized to carry the morphology; i.e., \textit{media tantum} forms occur.)\footnote{\citet{Inglese2022a} actually includes a third criterion:  The semantics [or function] of at least some verbs in (i) doesn’t match that of those in (ii), or vice versa. (In part, this excludes mono-functional morphemes.) The reader is referred to \citet{Inglese2022a} for further details; this criterion is not repeated in \citet{Inglese2023}.}
\end{enumerate}

\noindent As these features hold for the Ancient Greek middle inflection, Inglese states that Ancient Greek has a ``middle system.”\footnote{Fortuitously, we might say, for the history of linguistic terminology.} Citing \citet[417-425]{Delbruck1897} and \citet[431-432]{Romagno2010}, he notes that  \textit{-omai} \textsc{1sg.mid(dle)} occurs with verbs that alternate with non-middle \textit{-ō} \textsc{1sg.act(ive)} to mark valency change, as in \REF{ex:Payne:AGreekAlternate}; and it occurs with verbs that exclusively occur as middles, as in \REF{ex:Payne:AGreekMediatantum}. Valence-reducing and situational functions for the middle forms are in small caps in the right-most columns of these examples. Those in \REF{ex:Payne:AGreekAlternate} are found in \citet[490]{Inglese2022a}, and I have added those in  \REF{ex:Payne:AGreekMediatantum}.

%\newpage

\ea alternating \label{ex:Payne:AGreekAlternate}\\
\ea
\gll \textit{daí-ō}      \textit{daí-\textbf{omai}}   \textsc{anticausative}\\
  burn-\textsc{act}     burn-\textsc{mid}\\
\glt `burn (tr.)'     `burn (intr.)'\\

\ex
\gll \textit{kópt-ō}  	  \textit{kópt-\textbf{omai}}   \textsc{reflexive}\\
  hit-\textsc{act}     hit-\textsc{mid} \\
\glt    `hit (tr.)'     `hit oneself'\\

\ex
\gll \textit{dikáz-ō}     \textit{dikáz-\textbf{omai} }   \textsc{passive} \\
  judge-\textsc{act}     judge-\textsc{mid} \\
\glt    `judge (tr.)'     `be judged' \\
 \z 
 \z

 \ea non-alternating (\textit{media tantum}) \label{ex:Payne:AGreekMediatantum}\\
 \begin{tabular}{lll}
\textit{hê\textbf{mai}} & ‘sit’	& \textsc{body posture}\\
\textit{kínu\textbf{mai}} & ‘move (intr.)’ & \textsc{translational motion}\\
\textit{skúz\textbf{omai}} & ‘be angry at’ & \textsc{emotion}\\
\textit{mákh\textbf{omai}} & ‘fight’ & \textsc{reciprocal}\\
\end{tabular}
\z 

Kemmer discussed the \textsc{reflexive} function as perhaps the central function of a middle construction. Inglese finds \textsc{anticausative} is the most central function in a multi-dimensional mapping study, followed (in order) by \textsc{reflexive, passive, reciprocal, self-benefactive, antipassive, facilitative} and \textsc{impersonal}.  

Somewhat separately from thinking about middles only in terms of valency reducing functions, both Kemmer and Inglese discuss a range of what I call semantic ``situation types" that middles commonly express. Some of these are included in the list of functions (in small caps) for the examples just above and for the Italian \textit{si} construction illustrated in \REF{ex:Payne:Italian}.

%\newpage

\ea Italian \textit{si} construction, with situation types; adapted from \citet[492]{Inglese2022a}\label{ex:Payne:Italian}\\
\begin{tabular}{lll}
\textit{radersi} & ‘shave (self)’ & \textsc{grooming/reflexive}\\
\textit{alzarsi} & ‘stand up’ & \textsc{change in body posture}\\
\textit{girarsi}& ‘turn (intr.)’ & \textsc{non-translational motion}\\
\textit{spostarsi} & ‘move (intr.)’ & \textsc{translational motion}\\
\textit{arrabbiarsi} & ‘get angry’ & \textsc{emotion}\\
\textit{immaginarsi} & ‘envisage’ & \textsc{cognition}\\
\textit{sciogliersi} & `melt (intr.)'  & \textsc{spontaneous/anticausative}\\
\textit{combattersi} & ‘fight’ & \textsc{reciprocal}\\
\textit{colpirsi} & ‘hit oneself’ & \textsc{reflexive}\\
\textit{si va} & ‘one goes’ & \textsc{impersonal}\\
\textit{si vende} & ‘is sold’ & \textsc{passive}\\
\textit{si taglia (facilmente)} & ‘is (easily) cut’ & \textsc{facilitative}\\
\end{tabular}
\z 

With this brief background on approaches to defining middle constructions and systems, I now turn to the Maa data.

\section{Background on Maa}\label{sec:Payne:MaaBackground}

Maa is generally described as a verb initial language, though order of major constituents can vary. It is a mixed head- and dependent-marking language: verbs mark person and number of subjects and speech-act-participant objects, while post-verb nominal phrases and free pronouns display a marked nominative case system \citep{Konig2006}, expressed by tone melodies which vary by lexical item \citep{TuckerMpaayei1955}. Many morphemes have multiple allomorphs due to advanced or retracted tongue-root (ATR) vowel harmony, grammatical and lexical tone, and other morphophonemic processes.\footnote{In examples, an accute accent marks High tone, a macron marks Downstep High, a circumflex marks Falling, and Low tone is unmarked.} 

At a minimum, Maa verbs must index person (or number of subject for infinitives) via prefixes, and (roughly) finiteness distinctions via tone. But they can also be morphologically quite complex \citep[51-95]{TuckerMpaayei1955}. Due to three intersecting parameters, a single maximal verb template does not exist. Briefly, verb roots and stems divide into two morphological classes, known as Class 1 and Class 2 \citep[52]{TuckerMpaayei1955}, which differ in how causatives and perfect(ive) aspect are marked. Transitivity and stative versus dynamic lexical aspect (\textit{Aktionsart}) also affect what morphology a root or stem can take \citep{Payne2015b}.  For example: Already intransitive stems cannot further take middle or antipassive morphology; already stative stems do not normally take directionals (without shifting to a non-directional meaning), but can take the inchoative; dynamic stems can take directionals but not the inchoative.

Maa is rich in valence morphology. Valence-increasing morphology includes the dative applicative (expressing ‘goal-reached, benefactive, recipient’), the instrumental applicative (‘instrument, location, manner'), and causatives which differ for Class 1 versus Class 2.\footnote{The instrumental applicative doubles as the Class 2 causative, where an animate ``instrument" is a causee. Additionally, an external possession construction increases syntactic valence but does not involve any verb morphology \citep{Payne1997-extposs}.} Valence-decreasing constructions include the impersonal \citep{Greenberg1959, Payne2011}, antipassive \citep{Payne2021-apas}, and middles. These are illustrated where relevant in the paper.

Mood and grammatical aspect are salient categories. Morphological aspects include non-perfect(ive) (generally unglossed in examples), perfect(ive),\footnote{\citet{Konig1993} argues that the relevant forms are `perfective' in function; \citet{Payne2015a} shows they function in discourse like a `perfect'. Here I use the term ``perfect(ive)" to cover the range of functions discussed in those two works.} and progressive \citep{Konig1993, Payne2015a}.  The itive (`motion away') and ventive (`motion towards') directionals also have some aspectual, as well as other grammatical functions \citep{Payne2021-am}.

We will see that Maa middles interact in particular ways with certain other morphemes. For convenience, I refer to middle plus these other relevant items as ``core" affixes and in this paper simply stipulate them as the items in \tabref{tab:Payne:Core}. They include causative, antipassive, applicative, directional, and aspect morphology. All are suffixes except the Class 1 Causative.\footnote{“Core” morphology is probably old as affixal material. It mostly corresponds to what some might call “derivation” but it includes productive aspect morphology which some might consider inflectional. Participant number is expressed both in some core and non-core affixes. See \citet{Payne2015b} for fuller discussion.}

\begin{table}
\caption{Core affixes (major allomorphs)}
\label{tab:Payne:Core}
\begin{tabular}{ll}
\lsptoprule
All middle forms & (see Table \ref{tab:Payne:Middle})\\ 
\midrule
Progressive & \textit{-ɪta/-ito}\\
Perfect(ive) & (\textit{tV-})...\textit{-a(k)/-o(k)}\\
Antipassive & \textit{-ɪshɔ(r)/-isho(r)}\\
Itive directional & (see Table \ref{tab:Payne:Perfect-Itive})\\
Ventive directional & \textit{-ʊ(n)/-u(n)}\\
Dative applicative & \textit{-akɪ(n)/-oki(n)}\\
Instrument applicative & \textit{-ie(k)}\\
Class 1 Causative & \textit{ɪta-/ito-}\\
\lspbottomrule
\end{tabular}
\end{table}

%\newpage 

\tabref{tab:Payne:Middle} presents middle and middle-complex endings. Note that some alternate for perfect(ive) versus non-perfect(ive) aspect themselves. None of the middle affixes or affix complexes co-occur with progressive or antipassive affixes. This is likely for semantic reasons, as the progressive and most uses of the antipassive assume dynamic readings, contrasting with the general tendency for middle functions and senses to be non-dynamic (but notably, dynamicity is typical for reflexive and reciprocal functions of the middles). The middle endings can occur on directional, applicative, and causative stems, as partially seen in \tabref{tab:Payne:Middle}; or as we will see in the case of \textit{media tantum} verbs, they may completely replace them. Some middle forms vary for subject number and for aspect but it appears that applicative-plus-middle and directional-plus-middle combinations do not display number and aspect contrasts (aside perhaps from tone). 

\begin{table}
\caption{Middle morphology (major allomorphs)}
\label{tab:Payne:Middle}
\begin{tabular}{lll}
\lsptoprule
& Singular & Plural\\
\midrule
Non-Perfect(ive) Middle (\textsc{mid.npf}) & \textit{-a(r), -o(r), -ɔ(r)} & \textit{-ata, -oto, -ɔtɔ} \\
Perfect(ive) Middle (\textsc{mid.pf}) & \textit{-ɛ, -e} & \textit{-atɛ, -ote}\\
\\
Ventive+Middle & \multicolumn{2}{c}{\textit{-ʊnyɛ, -unie}}\\
Itive+Middle & \multicolumn{2}{c}{\textit{-ari, -ori}}\\
Dative+Middle & \multicolumn{2}{c}{\textit{-akino, -okino}}\\
Middle+Instrument & \multicolumn{2}{c}{\textit{-arɛ, -ore}}\\
Middle+Inchoative & \multicolumn{2}{c}{\textit{-ayu, -oyu}}\\
\lspbottomrule
\end{tabular}
\end{table}

In Maa middle constructions, the only expressable argument (which would be the patient or ``internal" argument of any corresponding transitive stem) is indexed as subject on the verb and occurs in the nominative case if expressed in post-verbal position by a lexical phrase or free pronoun. Compare the tonal case marking on `heart' in \REF{ex:Payne:active-trans}, where `heart' is in the absolute case (unglossed) as the transitive object, with its form in \REF{ex:Payne:active-trans2}--\REF{ex:Payne:middle}, where it has nominative case (glossed \textsc{nom}). In  \REF{ex:Payne:middle}, note that `heart' is the only argument of a middle verb.

\ea\label{ex:Payne:active-trans}
\gll k-ɛ́-ɛ́ta ɔl=\textbf{táʊ́} o-gól ɛlɛ́  tʊ́ŋání \\
\textsc{cn2}-3-have \textsc{msg}=heart \textsc{rel.msg}-be.hard that.\textsc{msg.nom} elder.\textsc{nom}\\
\glt `That man has a hard heart.' (e.g., doesn't show concern when s.o. dies) [\textit{ɔltáʊ́} = transitive object in absolute case]
\z 

\ea\label{ex:Payne:active-trans2}
\gll n-áa-lilîû ɔl=\textbf{tâʊ̂} naléŋ\\
\textsc{cn1}-\textsc{3>1sg}-stink \textsc{msg}=heart.\textsc{nom} very\\
\glt `I feel like vomiting.' (lit. `My heart stinks a lot.')
[\textit{ɔltâʊ̂} = transitive subject in nominative case]\footnote{Example \REF{ex:Payne:active-trans2} is an external possession construction in which the \textsc{1sg} object is interpreted as the possessor of the subject \citep{Payne1997-extposs}. Even more literally it would read `The heart stinks me very much.'}
\z 

\ea \label{ex:Payne:active-intrans}
\gll ɔl=tʊŋáni ɔ-wáŋ ɔl=\textbf{tâʊ̂} \\
\textsc{msg}=person \textsc{rel.msg}-be.bright \textsc{msg}=heart.\textsc{nom}\\ 
\glt `a person whose heart shines' (i.e., a person who has no bad intentions)
[\textit{ɔltâʊ̂} = intransitive subject of relativized verb, in nominative case]
\z 

\ea \label{ex:Payne:middle}
\gll n-é-dūŋ-ō \textbf{tâʊ̂}, n-é-yē \\
\textsc{cn1}-3-cut-\textsc{mid.npf} heart.\textsc{nom} \textsc{cn1}-3-die\\
\glt `(Her) heart stops (lit. is cut/cuts), and she dies.' (nochild.034)
[\textit{tâʊ̂} = argument of middle construction in nominative case] 
\z 

Finally, some middle allomorphs are quite similar phonologically to those of the perfect(ive) and itive. Thus, for clarity, \tabref{tab:Payne:Perfect-Itive} presents major allomorphs of the perfect(ive) and itive categories.

\begin{table}
\caption{Major perfect(ive) and itive allomorphs}
\label{tab:Payne:Perfect-Itive}
\begin{tabular}{ll}
\lsptoprule
Perfect(ive) & \textit{-a(k), -o(k)}\\
Itive & \textit{-aa, -oo, -ar(a), -or(o), -aya, -oyo}\\
\lspbottomrule
\end{tabular}
\end{table}

\section{Alternating Maa middles}\label{sec:Payne:Alternating}

I now illustrate alternating middle behavior. Middles of this type have non-middle transitive roots (or stems). Dynamic transitive roots quite productively take middle endings to express reflexive, reciprocal (in plural contexts), anticausative or spontaneous \citep{Haspelmath1987}, and result-state functions.\footnote{By \textsc{result-state}, I mean the state of an undergoer as the result of an action-process. The middle form itself expresses a state, while the transitive non-middle root would express an action-process potentially resulting in that state. Note that \textsc{result-state} is not intended to (necessarily) imply a passive reading where any sense of an \textsc{agent} is even semantically retained. Also, some Maa middle forms express a \textsc{state} where it is hard to see that state as the result of a transitive action-process.} Depending on semantics of the root, some alternating middles may also express movement, emotion, cognition, etc., as well as some quite idiosyncratic (hence, clearly lexicalized) meanings. Consider the range of senses that arise with middle forms of the dynamic transitive verb root \textit{duŋ} `cut' in \REF{ex:Payne:dung}; the listed senses for the intransitive middle forms are not necessarily exhaustive.

\ea \label{dung}\label{ex:Payne:dung}
\begin{xlist}
\ex 
\gll a-dúŋ		 \\
\textsc{inf.sg}-cut\\
\glt ‘to cut’ \hfill \textsc{dynamic, transitive}\\

\ex 
\gll a-duŋ-ó \\
\textsc{inf.sg}-cut-\textsc{mid.npf} \\
    \glt ‘to be in a cut state, demarcated, separated’  \hfill \textsc{result-state}\\
‘to cut self, divide self (into groups)’ \hfill \textsc{reflexive}\\
‘to cease/disrupted flow’ \hfill \textsc{anticausative, result-state}\\
 (of electricity, water, cows from heaven, heart beat, etc.)  \\
‘to be past menopause’ \hfill \textsc{spontaneous, state} \\
‘to be dried up’\footnote{Note that the transitive root \textit{duŋ} does not have a sense `dry sth.'; a completely unrelated transitive stem \textit{itoi} means `dry sth'.} \hfill \textsc{spontaneous, state}\\
 
 \ex 
 \gll áa-duŋ-o 		\\
 \textsc{inf.pl}-cut-\textsc{mid.npf}	\\
  \glt ‘to cut each other’ \hfill\textsc{reciprocal}\\
  ‘to be divided (into portions)’ \hfill \textsc{result-state}
  \end{xlist}
\z 

Additional examples of alternation between dynamic transitive and middle stems follow. In \REF{ex:Payne:predictable}, the middle has quite predictable senses. In \REF{ex:Payne:mentally-ill}, we see both predictable and metaphorical senses. In \REF{ex:Payne:unpredictable},  the middle has a rather unpredictable sense of `(spontaneously) coagulate', even though the operative valence decreasing function is a highly productive \textsc{anticausative} one.

\ea \label{ex:Payne:predictable}
\begin{xlist}
\ex 
\gll a-ɪnyɪnyɪrʊ́  \\
\textsc{inf.sg}-exude.liquid.\textsc{ven}\\
\glt ‘to let out tiny amounts of liquid' \hfill \textsc{dynamic/spontaneous,}\\
\- (e.g., sweat) \hfill \textsc{transitive}

\ex	
\gll a-ɪnyɪnyɪrʊnyɛ́ \\
\textsc{inf.sg}-exude.liquid.\textsc{ven.mid}\\
\glt ‘to be let out in tiny amounts’ 		 \hfill \textsc{anticausative, spontaneous}\\ 
\end{xlist}
\z 

\ea\label{ex:Payne:jing}
\begin{xlist}
\ex a-jɪ́ŋ\\
\textsc{inf.sg}-enter
\glt `to enter (a place)' \hfill \textsc{dynamic, transitive}

\ex\label{ex:Payne:mentally-ill}
\gll a-jɪŋ-á\\
\textsc{inf.sg}-enter-\textsc{mid.npf}\\
\glt `to be entered (of a place)' \hfill \textsc{result-state}\\
`to be possessed by evil' \hfill \textsc{(result-)state}\\
`to be(come) mentally ill' \hfill \textsc{cognition}\\
\end{xlist}
\z 

\ea \label{ex:Payne:unpredictable}
\begin{xlist}
\ex 
\gll a-ɪbʊ́ŋ \\
\textsc{inf.sg}-catch\\
\glt ‘to catch (sth. moving), arrest, touch'\\
`to rape, stick to following a path’ \hfill \textsc{dynamic, transitive}\\

\ex 
\gll a-ɪbʊŋ-á    \\
\textsc{inf.sg}-catch-\textsc{mid.npf}\\
\glt ‘to coagulate, be coagulated’ \hfill \textsc{spontaneous, state}

\end{xlist}
\z 

The preceding examples show that the various valence-reducing functions are not necessarily distinct when it comes to particular verbs: a single middle verb like \textit{a-duŋó} can have different functions like \textsc{reflexive} or \textsc{anticausative}. Also, given complex lexical semantics, a single middle stem can simultaneously express multiple situation types such as \textsc{motion}+\textsc{posture}, \textsc{emotion}+\textsc{cognition}, \textsc{emotion} plus some other aspect of behavior, etc.

Finally, in light of \citet{Inglese2022a}'s findings, one might ask whether Maa alternating middles do not express \textsc{passive}, \textsc{facilitative}, and \textsc{antipassive} valence-decreasing functions. In some contexts, the middle+inchoative form \textit{-ayu/-oyu} can express an `able to/will become' sense, which is close to, if not the same as, facilitative. This combination is seen in \REF{ex:Payne:unreturnable}--\REF{ex:Payne:mid-inch}.

\ea\label{ex:Payne:unreturnable}
\gll en=tóki		na-j-î				m-e-shúk-\textbf{ōy-ū}		\\
\textsc{fsg}=thing	\textsc{rel.fsg}-say-\textsc{imps}	\textsc{neg}-3-return.sth-\textsc{mid.npf}-\textsc{inch}	\\		
\glt `a thing that is/was called unreturnable' (renoi.038)
\z 

\ea\label{ex:Payne:mid-inch}
\gll n-é-pūō-ī		o-m-e-tú-duŋ-\textbf{oy-ú} ɔl=mʊ́rráni		ɛn=ámʊkɛ \\	
\textsc{cn1}-3-go.\textsc{pl-pl}	until-\textsc{pot}-3-\textsc{sbjv}-cut-\textsc{mid.npf-inch} \textsc{msg}=warrior	\textsc{fsg}=sandal.\textsc{nom}\\
\glt `They go until the warrior's shoe becomes cut (frayed, torn, separated).' (enamuke1.0002)\\
\z 
To my understanding, Maa does not really use the middle for the passive function because the impersonal typically expresses the idea that there is some unexpressed (and non-expressable) agentive ``people" that do the action, which is the passive sense \citep{Payne2011}; this is evident in \textit{najî} `that is called/which ``people" call' in \REF{ex:Payne:unreturnable} just above. The Maa middle -- in contrast to the impersonal -- completely erases the sense of an agent or causer, leaving just the undergoer. Finally, Maa has a dedicated antipassive suffix \textit{-ɪshɔ(r)/-isho(r)} \citep{Payne2021-apas}, which can be seen in \REF{ex:Payne:drunk-apas} below. The dedicated impersonal and antipassive forms carve out valence-reducing functions that the middle forms do not extend into (see also \citet{Inglese2022b} for some discussion).

\section{Maa-style \textit{media tantum}}\label{sec:Payne:MediaTantum}

In citation or their simplest person-inflected form, some Maa verbs always contain one or another of the middle endings in \tabref{tab:Payne:Middle}. Such verbs are reminiscent of Greek \textit{media tantum} verbs, but potentially differ in that they can alternate between the middle non-perfect(ive) (\textsc{mid.npf}) and middle perfect(ive) (\textsc{mid.pf}) aspect endings; or they drop the middle ending when they carry certain ``core" affixes. However, they cannot occur without a middle or some other core affix.

For example, consider \textit{a-mɛrá} `to be drunk/intoxicated, get (more) drunk'. \\
There is no synchronically recognizable verb *\textit{a-mɛ́r} without a final \textit{a}. It thus contrasts with alternating middle verbs described in the preceding section. 

Now, even though modern Maa has no simple verb \textit{*a-mɛ́r}, the final \textit{a} disappears in \REF{ex:Payne:drunk-apas}. The antipassive suffix \textit{-ɪ́shɔ̄(r)} presumably occurs on the bracketed stem \textit{ɪ́tɛ-mɛ́r} ‘cause to be drunk’, but the relevant point is that the antipassive immediately follows what looks like a non-middle root form \textit{mɛ́r}. 

\ea \label{ex:Payne:drunk-apas}
\gll n-é-ōk il=keék ɔɔ-[ɪ́tɛ-mɛ́r]-ɪ́shɔ̄ \\	
\textsc{cn1}-3-drink 	\textsc{mpl}=medicines \textsc{rel.mpl-caus}-be.drunk-\textsc{apas}\\
\glt ‘They take drugs that intoxicate.’ (malk01.1083)
\z 

The antipassive suffix begins with a vowel; so one might ask whether \emph{a} disappears from \textit{mɛra} due to phonological simplification of a vowel cluster \emph{aɪ}. The answer is ``no." Maa has no prohibition against vocalic clusters; at least 42 cluster types, including rising and falling diphthongs, are attested \citep{GriscomPayne2017}. Also, epenthetic [y] and [w] can break up certain vowel sequences. So deletion to simplify a vowel cluster is an unlikely explanation.

Most tellingly, a Class 1 causative prefix \emph{ɪta-}, alone, can result in dropping the final vowel from a \emph{media tantum} verb, demonstrated by \REF{ex:Payne:drunk-caus}. The contrasts in \REF{ex:Payne:pregnant} show the same patterns for \textit{a-nʊtá} `to be pregnant'. Similar alternation patterns occur for other Maa-style \textit{media tantum} verbs.\footnote{Syntactic repercussions vary in accord with what core affix “displaces” the historical middle (i.e., with  the top-most category of a  derived stem): an antipassive stem correlates with intransitive syntax, a causative stem with transitive syntax, etc. For reasons of space, I cannot deal much with derivational layers here.}

\ea\label{ex:Payne:drunk-caus}
\gll k=ɛ́-ɪ́tɛ́-mɛr\\
\textsc{cn2=3-caus}-be.drunk\\
\glt `S/he will make him/her drunk.'
\z 

\ea \label{ex:Payne:pregnant}
\begin{xlist}
\ex 
\gll a-nʊtá \\
\textsc{inf.sg}-be.pregnant\\
\glt ‘to be pregnant’
\ex *a-nʊ́t
\ex 
\gll a-ɪtʊ-nʊ́t 	\\
\textsc{inf.sg-caus}-be.pregnant \\	
\glt `to impregnate (a human)’ 
\end{xlist}
\z 

For \textit{a-mɛrá}, historical evidence suggests there may have been a root \emph{*mɛr} at some point in Nilotic history. \citet[357]{Vossen1982} reconstructs \emph{*-mɛr-a} (parsing off the final \emph{*a}) for Proto-Teso-Lotuko-Maa, as his data shows all cognates within this sub-branch of Eastern Nilotic as having a middle-like ending or suffix. Notably, he presents the Ongamo form
\emph{a-tɛ́-mɛ́r-ɛ̀}, which would correspond exactly to a Maa \textsc{mid.pf} form of a Class 1 (but synchronically non-existing) root \emph{*mɛr}. The alternation between \textit{a} and \textit{ɛ} on Eastern Nilotic cognates of ‘be drunk, intoxicated’ argues that the final \textit{a} in Maa \textit{a-mɛrá} corresponds (historically) to the \textsc{mid.npf} affix \textit{-a(r)}. That is, \textit{mɛra} is not a historically basic stative root which just coincidentally happens to end in the phoneme \textit{a}.

In sum, there is evidence that Maa-style \textit{media tantum} verbs contain a now-lexicalized historical \textit{-a/-o} middle ending. This ending is part and parcel of such a verb in its most basic form, such that native speakers recognize no meaning without it. But the fact that it disappears when certain other core affixes occur is a tell-tale sign that the \textit{a}\sim\textit{o} is the historical middle ending. (In \sectref{sec:Payne:Spurious} we will see that ending in an  \textit{a}\sim\textit{o} vowel, even if having the ``right" kind of semantics, does not necessarily mean a verb belongs to the \textit{media tantum} set.)

To give a better sense of the range of semantics found with Maa-style \textit{media tantum} verbs, \tabref{tab:Payne:MT} lists a few additional such verbs with their situation types; \sectref{sec:Payne:Motivation} will further address semantic issues.

\begin{table}
\caption{Selected \textit{media tantum} verbs and situation types}
\label{tab:Payne:MT}
\fittable{
\begin{tabular}{lll}
\lsptoprule
a-rri̩á̩	& `to fall down' & \textsc{translational movement} \\
a-isiijó & `to ferment' & \textsc{spontaneous}\\
a-ɪbalá & `to be conspicuous; unambiguous; \\
& likely to happen; true' & \textsc{state, facultative} \\
a-ineníá & `to be piled up, crammed' & \textsc{(result) state}\\
a-lioó & `to be visible, apparent' & \textsc{state}\\
a-lubó & `to be underfed (cattle)' & \textsc{state}\\
a-narikinó & `to match; be advisable' & \textsc{state, facultative}\\
a-ɪkɪrɪkɪrá & `to quake (earth)' & \textsc{non-human process}\\
a-sɛshá & `to be sickly' & \textsc{body-state}\\
a-ɪgʊtʊmá & `to squat, sit' & \textsc{body-posture}\\
a-gʊrrʊmá  & `to vomit' & \textsc{body-process}\\
a-dalá & `to be inattentive; play joyfully' & \textsc{emotion/cognition,}\\
& & \textsc{human behavior}\\
a-ɪmasó & `to be proud' & \textsc{emotion/cognition}\\
a-ŋɪdá & `to be happy; haughty; naughty' & \textsc{emotion},\\ 
& & \textsc{human propensity,}\\
& & \textsc{human behavior}\\
\lspbottomrule
\end{tabular}
}
\end{table}

\section{Partially assimilated middles}\label{sec:Payne:PartialAssim}
\largerpage[2]
I now turn to verbs which appear to be partially assimilated into the \textit{media tantum} group. \textit{A-ɪgará} ‘to obscure, hide sth. behind sth.’ behaves like a \textit{media tantum} verb in that the form \textit{*a-ɪgár} does not occur, but the final \textit{a} is dropped when other core affixes are added. Unlike the vast majority of alternating middle and \textit{media tantum} stems, \textit{a-ɪgará} can function in a transitive sentence frame, even with the final \textit{a}, as in \REF{ex:Payne:obstruct}.\footnote{Such mismatch between intransitive morphological form and transitive syntactic behavior is known as ``deponency" \citep{Grestenberger2023}. Aside from \textit{a-ɪgará}, I do not know of deponent behavior among other Maa \textit{media tantum} verbs.}

\ea \label{ex:Payne:obstruct}
\gll ɛ-ɪgára ɛnk=áji en=kíné  \\
3-obstruct \textsc{fsg}=house.\textsc{nom} \textsc{fsg}=goat\\
\glt `The house will obstruct/obstructs the goat (from view).'
\z 

Like alternating middles, \textit{a-ɪgará} can have a reflexive sense, as in \REF{ex:Payne:obstruct-refl}. This suggests the final \textit{a} might correspond to the historical middle suffix (note the tone difference in the verbs in \REF{ex:Payne:obstruct} versus \REF{ex:Payne:obstruct-refl}). And like \textit{media tantum} verbs, final \textit{a} is dropped if a core affix is added. We see this in \REF{ex:Payne:obstruct-caus} when the Class 2 causative/instrumental is added. Further, the progressive is normally disallowed with a stative base, but it occurs in \REF{ex:Payne:obstruct-prog} directly on the form \textit{ɪgar}.\footnote{This example is from a speaker of IlKeekonyokie Maa. Another speaker of Purko Maa states that the verb form makes sense though he ``hasn't used it before".} 

\ea \label{ex:Payne:obstruct-refl}
\gll  	ɛ-ɪgárā		ɔl=páyian (ɔl=chaní) ŋolé \\	
3-hide.\textsc{pf}	\textsc{msg}=elder.\textsc{nom} \textsc{msg}=tree yesterday \\
\glt `The man hid himself (behind a tree) yesterday.'\\
\z

\ea \label{ex:Payne:obstruct-caus}
\gll	a-ɪgar-íé		em=búku	ol=órika 	\\
\textsc{inf.sg}-hide-\textsc{caus}	\textsc{fsg}=book	\textsc{msg}=chair	\\
\glt ‘to hide a book behind a chair’
\z

\ea \label{ex:Payne:obstruct-prog}
\gll ɛ-ɪ́gar-ɪ́ta	ɛnk=áji	en=kíné  \\
 3-hide-\textsc{prog}	\textsc{fsg}=house.\textsc{nom}	\textsc{fsg}=goat\\
 \glt ‘The house is obstructing the goat (from view).’
 \z 

These facts suggest that \textit{ɪgar} was (historically) a dynamic root and the final \textit{a} might be the historical middle, allowing either a reflexive/reciprocal sense or the rather stative (though transitive) notion of ‘obstruct from view, block’. (Subjects of this verb are easily inanimate, also potentially correlating with a less dynamic profile.) More specifically, these facts might suggest that \textit{a-ɪgará} is a Maa-style \textit{media tantum} verb. Where it notably differs from that group of verbs, however, is that it does not take the \textsc{mid.pf} suffix \textit{-ɛ}. Rather, its \textsc{pf} counterpart varies tonally from the non-\textsc{pf} form. Thus, the \textsc{pf} form looks like a Class 2 non-middle verb. Compare \REF{ex:Payne:obstruct} above with \REF{ex:Payne:obstructed}:

\ea \label{ex:Payne:obstructed}
\gll ɛ-ɪgárā	ɛnk=áji	en=kíné 	\\
3-hide.\textsc{pf}	\textsc{fsg}=house.\textsc{nom}	\textsc{fsg}=goat 	\\
\glt ‘The house obstructed the goat from view.’
\z

The \textsc{pf} form suggests the hypothesis that there was a regular (historical) transitive root \textit{*ɪgar}, which still uses the productive Class 2 perfect(ive) \textit{-a(k)}, surfacing here with Downstep High tone (marked by the macron). In other respects, \textit{a-ɪgará} appears to be part of the \textit{media tantum} group.

As another type of partially assimilated \textit{media tantum} verb, consider \textit{a-sɪnyá} `to be holy, perfect, spotless', seen in a relativized form in \REF{ex:Payne:holy}.

\ea\label{ex:Payne:holy}
\gll ɔl=kɪtɛ́ŋ	ɔ-sɪ́nya	oshî ɛ-pɪk-ɪ́	in=taléŋo \\
\textsc{msg}=bovine	\textsc{rel.msg}-be.holy	always	3-put-\textsc{imps}	\textsc{fpl}=ceremonies \\
\glt `An ox which has no physical fault is put (slaughtered) in ceremonies.' [pure black or white, not spotted, nor an ``unsaintly” colour]
\z 

\citet[372]{Mol1996} states this verb occurs in the ``present tense” (i.e., non-\textsc{pf} aspect) only. That is, the final \textit{a} does not alternate with the \textsc{mid.pf} ending \textit{-ɛ}. In other ways, the base behaves as a Maa-type \textit{media tantum} verb: The sequence \textit{sɪny} is immediately followed by the ventive+middle form \textit{-ʊnyɛ} in \REF{ex:Payne:blessed}; and there is no final \textit{a} in the causative form in \REF{ex:Payne:caus-holy}.

\ea \label{ex:Payne:blessed}
\gll A-sɪny-ʊnyɛ́		naá	doí		e=síáâî ná-tií	ɛnk=arruoísho	atúā?\\
3-be.holy-\textsc{ven.mid} 	\textsc{focus}	indeed \textsc{fsg}=work.\textsc{nom} \textsc{rel.fsg.nom}-be.at	\textsc{fsg}=conflict	inside \\
\glt ‘Can work grow blessedly which has conflict (chaos) in it?'
\z 

\ea \label{ex:Payne:caus-holy}
\gll a-ɪtɪ-sɪ́ny \\
\textsc{inf.sg-caus}-be.holy \\
\glt ‘to make holy, cleanse’
\z 

\textit{A-sɪnyá} is clearly related (though in some as-yet unclear diachronic way) to the adjectives \textit{sinyáti}/\textit{sinyát}  `holy, blameless (\textsc{sg}/\textsc{pl})’. This, plus the fact that there is no perfect(ive) form -- neither \textit{-ɛ }nor a \textsc{pf} tonal form like in \REF{ex:Payne:obstruct-refl} and \REF{ex:Payne:obstructed} -- suggests that the diachronic route for how \textit{a-sɪnyá} may have entered or be entering the \textit{media tantum} group may differ from that of \textit{a-ɪgará}. Conceivably, however, its partial media tantum behavior is motivated by semantic similarity to other \textit{media tantum} verbs (an issue discussed in \sectref{sec:Payne:Motivation}).

\section{Spurious ``look alikes"}\label{sec:Payne:Spurious}

Sections \ref{sec:Payne:MediaTantum} and \ref{sec:Payne:PartialAssim} have argued that various stative citation-form verbs ending in \textit{a} or \textit{o} are \textit{media tantum} verbs. However, not all stative citation-form verbs ending phonologically in this way belong to this group. A case in point is \textit{a-ɪrowúá} `to be hot’. Though semantically stative, it is not a \textit{media tantum} verb. First, there is no alternation with \textit{-ɛ} for \textsc{pf} aspect; instead, it takes the regular Class 2 perfect(ive) suffix \textit{-a(k)}, as in \REF{ex:Payne:hot}. In this respect, it might seem like \textit{a-igará} `to obscure', discussed above. But secondly, the relevant \emph{a} remains when various core affixes are added, as shown in \REF{ex:Payne:hot-inst}. Indeed, various suffixes show that the root has a final (historical) \textit{dʒ} (orthographic \textit{j}) -- rather than a final \textit{r} which would be expected if it carried the (historical) non-\textsc{pf} middle.

\ea \label{ex:Payne:hot}
\gll ɛ-ɪrɔ́wūāj-ā táatá \\
3-be.hot-\textsc{pf}	today \\
\glt `It became hot today.’ / `It was hot today.’ / `It has been hot today.’
\z 

\ea \label{ex:Payne:hot-inst}
\gll a-irowuaj-íé \\
\textsc{inf.sg}-be.hot-\textsc{inst/caus}\\
\glt `to heat with' / `to make sth. hot'
\z 

Also, the inchoative suffix can be added, as is true generally for stative verbs, but not the middle+inchoative form \textit{-ayu/-oyu}:

\ea\label{ex:Payne:become-hot}
\gll a-irowuaj-ú\\
\textsc{inf.sg}-be.hot-\textsc{inch}\\
\glt `to become hot'
\z 

Some vowel-final stative verbs can be shown to end in a historical \textit{r} but nevertheless do not pattern as \textit{media tantum} forms. Consider \textit{a-adɔ́} `to be linearly extended' in \REF{ex:Payne:be.long}. The (historical) \textit{r} is shown by the inchoative derivation -- making one think the otherwise final \textit{ɔ} might be a middle allomorph. However, it is neither a middle nor \textit{media tantum} verb: Its final vowel does not disappear in the causative, or with any other core affix, nor does it alternate with \textit{ɛ} in \textsc{pf} aspect.

\ea \label{ex:Payne:be.long}
\begin{xlist}
\ex 
\gll a-adɔ́\\
\textsc{inf.sg}-be.linearly.extended\\
\glt `to be long/tall'

\ex
\gll ɛ-adór-ū ɔl=chánī\\
3-be.linearly.extended-\textsc{inch} \textsc{msg}=tree\\
\glt `The tree will become tall.'

\ex 
\gll í-ntɔ́-ɔ́dɔ 	\\
2-\textsc{caus}-be.linearly.extended\\		
\glt ‘You make him/her/it be tall/long.’
\end{xlist}
\z 

Another spurious case is \textit{a-naná} `to be tender, soft, gentle' which semantically would appear to be a prime middle candidate. However, it lacks a perfect(ive) counterpart with \textit{-ɛ}, distinguishing it from regular middle/\textit{media tantum} verbs \citep{Mol1996}, and it keeps final \textit{a} in conditions that trigger loss in \textit{media tantum} verbs; note the causative in \REF{ex:Payne:cause-be.soft}.

\ea \label{ex:Payne:cause-be.soft1}
\gll in=kírí	naá-nana	\\
\textsc{fpl}=meats	\textsc{rel.fpl}-be.soft 	\\
\glt `meats that are soft’ (inkiri.057)
\z 

\ea \label{ex:Payne:cause-be.soft}
\gll k=é-lō a-ɪta-naná \\
\textsc{cn2}=3-go \textsc{inf.sg}-\textsc{caus}-be.soft \\
\glt `S/he will make it (e.g., a hide) soft.'
\z

Unlike with `to be hot', no final consonant surfaces in the inchoative derivation of `be soft'; compare \REF{ex:Payne:become-soft} with \REF{ex:Payne:become-hot} above.

\ea \label{ex:Payne:become-soft}
\gll a-naná-ú \\
\textsc{inf.sg}-be.soft-\textsc{inch}\\
\glt `to become soft, etc.’
\z 

\section{What motivates development of \textit{media tantum} forms?}\label{sec:Payne:Motivation}

To explore what might motivate the development of \textit{media tantum} verbs, I tabulated senses and situation types of about 450 alternating middle and Maa \textit{media tantum} verbs.\footnote{These were middle forms that happened to be in a lexicography database, with well over 2750 senses involving verb forms. Since middle formation from active transitive verbs is highly productive, many potential alternating middles happen to not be included in the database.} The specific situation types for different senses are listed in \tabref{tab:Payne:SitType1}. 

%Real table
\begin{table}
\caption{Situation types evaluated for Maa middle verbs}
\label{tab:Payne:SitType1}
\begin{tabular}{ll}
\lsptoprule
Body-state & Body-posture  \\
Body-process & Non-human process \\
Emotion/Cognition & Facilitative (most with inchoative) \\
Human behavior  &Human propensity \\
 (e.g., once) &(typical characteristic)\\
Reflexive (grooming, action, etc.) & Reciprocal\\
Motion-translational  (most with directional)& Motion-nontranslational\\
Result-State  & State (other; not from process)  \\
Spontaneous  \\
\lspbottomrule
\end{tabular}
\end{table}

To appreciate how senses were categorized for situation type, \tabref{tab:Payne:SitTypes} presents a few middle stems, their senses, and corresponding situation types. This was done for all 450 middle stems in the data. (For reasons of space, the middle stems in Table \ref{tab:Payne:SitTypes} are not full word forms, as they lack an infinitive or person prefix.)

\begin{table}[t]
\caption{Middle verb situation types and senses  (selected examples and categories)}
\label{tab:Payne:SitTypes}
\begin{tabular}
{
>{\RaggedRight\arraybackslash}p{1.5cm}
>{\RaggedRight\arraybackslash}p{2cm}
>{\RaggedRight\arraybackslash}p{1.9cm}
>{\RaggedRight\arraybackslash}p{1.9cm}
>{\RaggedRight\arraybackslash}p{2cm}
}
\lsptoprule
middle stem  & \textsc{body-process} & \textsc{emotion/ cognition} & \textsc{human propensity} & \textsc{human behavior} \\
\lsptoprule
\textit{raya} & & emotionally cling to, have exclusive relationship with & love to be s.where or have sth. against contrary forces & elope (of a girl)\\
\midrule
\textit{ŋɪdarɛ}  & & be proud of & & boast of\\
\midrule
\textit{peepeeno} & & & be trouble-maker & loiter \\
\midrule
\textit{wuapa} & & & be insane, abnormal due to drugs & overreact angrily, ``snap" \\
\midrule
\textit{oko} & & & & sing about exploits\\
\midrule
\textit{gɪra} & do slowly, listlessly & & be naturally quiet, soft-spoken & be silent, quiet\\
\midrule
\textit{liyio} & & be lonely\\
\lspbottomrule
\end{tabular}
\end{table}

Frequencies of situation types relative to alternating and \textit{media tantum} stems were then examined. Eventually, I grouped those situation types judged to be more typically human-related (HR), versus those typically less human-related (``other"; \textsc{oth}). The HR types were \textsc{body-state, body-posture, body-process, cognition/emotion, human propensity}, and \textsc{human behavior}. The \textsc{oth} types were \textsc{non-translational motion, spontaneous, state, result-state, facilitative}, and \textsc{non-human process}. Reflexive/reciprocal and translational motion were now excluded from further statistical evaluation, as the former is so highly productive with nearly all dynamic transitive stems and the latter is so highly productive with derived directional stems. The reflexive/reciprocal functions are dominantly human-related,\footnote{Of the 28 reflexive functions that happened to be included in the senses of 450 middle stems, only four or five had grooming-type senses and all pertained to alternating middles: \textit{a-itukúó} `to bathe self', \textit{a-ɪsʊjá}  `to bathe, wash self/hands', \textit{a-ɪsɪkarrá} `to adorn self',\textit{ a-ishopó} `to dress self', and if one considers it to be a grooming event also \textit{a-girro}́ `to scratch self (e.g., due to itching)'.} and the middles of translational motion are dominantly non-human related, and I wanted to evaluate any potential skewing among the other situation types. This left 402 middle stems.

The distribution of situation-type tokens in the sample of 402 middle stems is given in \tabref{tab:Payne:Stats}. Since a given verb stem can have more than one sense and instantiate more than one situation type, each token of a situation type is reflected in the table: 313 total for alternating middle stems, and 113 total for \textit{media tantum} stems (N=426). That is, the N's in the table do not represent numbers of middle verb stems, but rather instances of situation types. The difference in how often \textit{media tantum} forms correlate with typically human-related situations versus ``other" situations is statistically significant (χ2 = .00365, p < .01). 

%Table to modify figures in
\begin{table}
\caption{Typically human-related (HR) and other (OTH) situation types relative to alternating versus media tantum stems}
\label{tab:Payne:Stats}
\begin{tabular}{lrr}
\lsptoprule
\textsc{situation} & \textsc{alternating} & \textsc{media tantum} \\
& N=313 & N=113 \\
\midrule
Less typically human-related (OTH)\\
Motion-nontranslational & 18 & 1 \\
Spontaneous & 18 & 4 \\
State/Result-State & 111 & 26 \\
Facilitative & 8 & 2 \\
Non-human process & 13 & 2 \\
(Subtotals) & (168) & (35) \\
\midrule
Typically human-related (HR)\\
Body-state & 25 & 12 \\
Body-posture & 13 & 15 \\
Body-process & 26 & 16 \\
Emotion/Cognition & 30 & 15 \\
Human propensity & 30 & 15 \\
Human behavior  & 21 & 9 \\
(Subtotals) & (145) & (78) \\
\lspbottomrule
\end{tabular}
\end{table}


\begin{figure}
\caption{``Other" (OTH) versus typically ``Human Related" (HR) Situation Types in Alternating (Alt) and \textit{media tantum} (MT) Verbs}
\label{fig:Payne:1}
\begin{tikzpicture}
\begin{axis}
[
ybar,
nodes near coords,bar width=0.5cm,
symbolic x coords={ ,OTH,HR, },
enlarge x limits=1.5,
  legend style={
    at={(1, .6)},
    anchor=south
    }
]
\addplot[style={gray,fill=gray}] coordinates{(,) (OTH,168)(HR,145)};
\addplot[style={black,fill=black}] coordinates{(,)(OTH,35)(HR,78)};
\legend{Alt gray, MT black}
\end{axis}
\end{tikzpicture}
\end{figure}

It may be hard to see the tendencies from \tabref{tab:Payne:Stats}.  Thus, \figref{fig:Payne:1} depicts the ``other" (OTH) versus typically human-related (HR) situation types relative to alternating (Alt) versus \textit{media tantum} (MT) types. The figure shows that a higher percentage of Maa \textit{media tantum} verbs correlate with human-related situations, compared to the percentage of all alternating middles that correlate with human-related situations. This pattern accords with \citet{Lyons1968}'s general characterization of what middle verb forms are about, namely, that they express concepts that concern or affect the interests of the subject of the verb. (But perhaps contra Lyons, the pattern in \figref{fig:Payne:1} does not support the hypothesis that \textbf{all} middles have this semantic profile, as the percentage of alternating verbs with ``other" functions outweighs the percentage of human-related functions.)

\largerpage
Quite clearly, the \textit{media tantum} verbs do not have a preponderance of reflexive uses; so at least the Maa \textit{media tantum} group seems to generally counter \citet{Kemmer1993}'s characterization of the middle prototype as involving the reflexive function. However, reflexive, reciprocal, and anticausative (potentially leading to \textsc{result-state}) functions are productive for alternating middles.\footnote{For Acholi (Western Nilotic), \citet[194--196]{Kemmer1993} traces a middle marker \textit{-ɛ̂} to a reflexive ending which she represents as *-\textit{ri} (but without actual reconstruction). To my knowledge, no one has reconstructed a middle suffix for Eastern Nilotic, but grammars of two Eastern Nilotic languages closely related to Maa, Lopit and Ateso, note middle \citep{Moodie2019} and reflexive \citep{Barasa2017} suffixes, respectively, with the form \textit{-a/-o}. See \citet{Inglese2023} for a broader discussion of the  diachronic typology of middle markers.}

The pattern in \figref{fig:Payne:1} leads to the hypothesis that higher usage frequency of middle forms expressing typical human situations may promote their lexicalization in such situations, followed by ultimate loss of the dynamic transitive counterparts. For some comparison, \citet{ZombolouAlexiadou2014} find that in a corpus study of non-alternating middle-form verbs in Modern Greek, a majority (89\%) have an ``experiencier" meaning \citep[337]{ZombolouAlexiadou2014}; and an ``experiencer" is typically human. Of course, the semantics of ``human situations" does not explain all of the \textit{media tantum} instances in Maa; but it is likely the explanation for lexicalization of forms like \textit{a-mɛrá} `to be drunk, intoxicated'.

\section{Conclusions}\label{sec:Payne:Conclusions}

This study concurs that Maa has a middle system in the sense of \citet{Inglese2022a} with alternating middle and \textit{media tantum} verbs, both with a range of functions. It goes beyond previous work by demonstrating the specific nature of ``Maa-type \textit{media tantum}" verbs and showing that some forms share only certain \textit{media tantum} features. The analysis thus is also a case study of a  linguistic category with permeable boundaries. A category consists of a web of signs (constructions), each member with its own semantic profile and frequency of use. New items can get (partially) pulled into a category based on perceived similarity in form and meaning. Alternatively, items can become marginalized in a category, lose full productivity, and may eventually disappear from a category. This may be happening with verbs like \textit{a-ɪgará} `to obscure', which appear to be partially assimilated to the \textit{media tantum} category, as their simple root forms have lost full productivity as members of the alternating middle set.

\section*{Abbreviations}
\begin{tabularx} {.45\textwidth}{>{\scshape}lQ}
\textsc{act} & active\\
\textsc{apas} & antipassive\\
\textsc{caus} & causative\\
\textsc{cn1} & connective 1 (high continuity)\\
\textsc{cn2} & connective 2 \\
\textsc{fpl} & feminine plural\\
\textsc{fsg} & feminine singular\\
\textsc{imps} & impersonal\\
\textsc{inch} & inchoative\\
\textsc{inf} & infinitive\\
\textsc{inst} & instrument\\
\normalfont{intr.} & intransitive\\
\textsc{mid} & middle\\
\end{tabularx}
\begin{tabularx} {.45\textwidth}{>{\scshape}lQ}
\textsc{mpl} & masculine plural\\
\textsc{msg} & masculine singular\\
\textsc{neg} & negative\\
\textsc{nom} & nominative case\\
\textsc{npf} & non-perfect(ive)\\
\textsc{pf} & perfect(ive)\\
\textsc{pl} & plural \\
\textsc{pot} & potential\\
\textsc{prog} & progressive\\
\textsc{rel} & relative (clause)\\
\textsc{sbjv} & subjunctive\\
\textsc{sg} & singular\\
\normalfont{tr.} & transitive\\
\textsc{ven} & ventive\\
\end{tabularx}


\section*{Acknowledgements}
I am grateful for the help of many Maa speakers, including Leonard Ole Kotikash, Sarah Tukuoo, Keswe Ole Mapena, Daniel Nalangu, Stephen Ole Muntet, †Neema Labani, Lazaro Ole Melubo, †Peter Ndetio, and Loatha Lesilale. This research has been partially supported by NSF grant SBR-9616482, the Fulbright Foundation, Mkwawa Campus of the University of Dar es Salaam, University of Nairobi, Africa International University, the University of Oregon, and SIL Global. Sincere thanks to all.

\sloppy %this command eliminates a quirk that appears in the bibliography, you can just leave it here.
\printbibliography[heading=subbibliography,notkeyword=this]

\end{document}

